%% RoadScribe -- manuscript
%%
%% Build:   latexmk -pdf main.tex      (or: pdflatex, bibtex, pdflatex, pdflatex)
%%
%% Pick ONE document class line:
%%   Single-column draft (default) ... \documentclass[manuscript,screen,review,anonymous]{acmart}
%%   MobiSys / two-column ............ \documentclass[sigconf,screen,review,anonymous]{acmart}
%%   IMWUT (PACM journal) ............ \documentclass[acmsmall,screen,review,anonymous]{acmart}
%% For the camera-ready version, drop "review,anonymous" and fill in the rights block below.
\documentclass[manuscript,screen,review,anonymous]{acmart}

\usepackage{algorithm}
% Lets TeX stretch inter-word space a little more before overfilling a line
% (two-column layouts with long technical terms).
\setlength{\emergencystretch}{1.5em}
\usepackage{algpseudocode}

% ------------------------------------------------------------------ macros
\newcommand{\R}{\mathbb{R}}
\newcommand{\ind}{\mathbf{1}}
% Red marker for every number that is not yet measured. The draft must not
% leave this state with any \TBD remaining.
\newcommand{\TBD}[1][]{\textcolor{red}{[TBD#1]}}

% ------------------------------------------------------------------ rights block (fill in when accepted)
\setcopyright{none}
\settopmatter{printacmref=false}
\renewcommand\footnotetextcopyrightpermission[1]{}
% Neutral running header for the draft; replace with the venue lines below.
\acmConference[Draft]{RoadScribe manuscript draft}{November 2026}{}
\acmYear{2026}
\acmDOI{}
\acmISBN{}
% IMWUT example:
% \acmJournal{IMWUT}\acmVolume{10}\acmNumber{1}\acmArticle{1}\acmYear{2027}\acmMonth{3}
% MobiSys example:
% \acmConference[MobiSys '27]{The 25th Annual International Conference on Mobile Systems, Applications and Services}{June 2027}{TBD}

\begin{document}

\title[RoadScribe]{RoadScribe: Grounded Trip Narratives from Phone Vibration and Map Topology with an On-Device Language Model}

\author{Anonymous Author(s)}
\affiliation{%
  \institution{Anonymous Institution}
  \country{Anonymous}}

\begin{abstract}
Phones in cars measure how roads shake the vehicle, but road-monitoring systems reduce this to a label or a roughness index per segment, and language models that read sensor data state values that nobody checks.
RoadScribe turns a trip recorded by a phone into a short, map-anchored narrative and answers follow-up questions, with every number, place and event traceable to a measured fact.
Layer~1 estimates the phone's mount per handling episode with a gated, rotation-equivariant gravity estimate and a closed-form Procrustes heading fit, applies the ISO~2631-1 weightings exactly, separates how the ride felt from how the road is, and labels impulsive events with an evidence-based maneuver gate.
Layer~2 segments the trip with Map-Informed PELT, whose penalty drops at OpenStreetMap transitions; we prove the pruning rule that keeps it exact under such penalties, where the textbook rule returned sub-optimal partitions in 136 of 300 random series.
Layer~3 compiles a typed claim ledger and a token-budgeted capsule, lets an on-device small language model plan over ledger identifiers under a grammar, fills in every value deterministically, and verifies each sentence, falling back to a template narrator when verification fails; on controlled errors in template narratives the verifier rejected 98.9\%.
On the PVS dataset the heading estimate recovered a rigid mount to within about $2^\circ$ across trips and exposed a $14^\circ$ mounting offset in one vehicle.
We define a benchmark whose metrics separate language errors from sensing errors, report it for three narrators on PVS, and specify a preregistered within-subject study with 28 drivers.
\TBD[: headline benchmark and study results.]
\end{abstract}

\begin{CCSXML}
<ccs2012>
 <concept>
  <concept_id>10003120.10003138.10003141</concept_id>
  <concept_desc>Human-centered computing~Ubiquitous and mobile computing systems and tools</concept_desc>
  <concept_significance>500</concept_significance>
 </concept>
 <concept>
  <concept_id>10010147.10010178.10010179.10010182</concept_id>
  <concept_desc>Computing methodologies~Natural language generation</concept_desc>
  <concept_significance>500</concept_significance>
 </concept>
 <concept>
  <concept_id>10010405.10010481.10010485</concept_id>
  <concept_desc>Applied computing~Transportation</concept_desc>
  <concept_significance>300</concept_significance>
 </concept>
</ccs2012>
\end{CCSXML}
\ccsdesc[500]{Human-centered computing~Ubiquitous and mobile computing systems and tools}
\ccsdesc[500]{Computing methodologies~Natural language generation}
\ccsdesc[300]{Applied computing~Transportation}

\keywords{mobile sensing, road surface, whole-body vibration, ISO 2631-1, change-point detection, data-to-text, grounded generation, small language models, on-device AI}

\maketitle

% =====================================================================
% Section 1 -- Introduction
% Usage: % =====================================================================
% Section 1 -- Introduction
% Usage: % =====================================================================
% Section 1 -- Introduction
% Usage: \input{sec1_intro} inside an acmart document (see main.tex).
% =====================================================================
\providecommand{\TBD}[1][]{\textcolor{red}{[TBD#1]}}

\section{Introduction}
\label{sec:intro}

Every smartphone that rides in a car carries an accelerometer, a gyroscope and a GNSS receiver.
Two decades of mobile sensing have used them to monitor roads: early systems detected potholes and braking from phones and dedicated in-vehicle units~\cite{mohan2008nericell,eriksson2008pothole}, and a large body of later work classifies surface types or regresses a roughness index from vehicle vibration~\cite{menegazzo2021pvs}.
The output of these systems is a label or a number per window or per road segment.
The people who drive the roads ask something different.
A commuter wants to know which part of the route was rough and whether to drive it differently.
A passenger prone to car sickness wants to know which stretches and which driving made the ride unpleasant.
A fleet manager wants to know where the vans take shocks, and whether the map is wrong about a surface.
These are questions about a specific trip, in language, and their answers combine facts that a single number cannot hold.

\paragraph{Why a scalar roughness score is the wrong output.}
The International Roughness Index (IRI) is defined on a measured road profile, through a reference quarter-car model driven at a fixed speed~\cite{sayers1986iri}.
A phone in a holder observes something else: the response of an unknown vehicle, through an unknown mount, at whatever speed the driver chose.
Regressing IRI, or any single roughness score, from such data runs into four problems.
\begin{enumerate}
\item \emph{It confounds the road with the drive.} Vibration grows with speed, so a regressor can score well by learning speed. On the PVS dataset~\cite{menegazzo2021pvs}, speed alone separates the three surface types with an AUC of 0.977 \TBD[: re-run audit], and an earlier ``kinetic normalization'' of our own encoded speed almost perfectly (Spearman 0.995 with speed) \TBD[: re-run audit].
\item \emph{It collapses distinct facts.} How the ride \emph{felt} depends on speed; how the \emph{road} is should not. A stretch with one violent speed bump and a stretch of steady cobblestone can share a mean vibration level, yet ISO~2631-1 treats the first as shock-dominated and requires different measures for it~\cite{iso2631}.
\item \emph{It has no place, cause or confidence.} A number does not say where along the street the roughness was, whether a jolt came from the road or from pulling away, or how sure the system is.
\item \emph{It cannot say what is missing.} When the GNSS fix drops out, the gyroscope is absent or the map contradicts the signal, a score is silently wrong; a sentence can say so.
\end{enumerate}

\paragraph{Why not a language model on the raw signal, or a template?}
Large language models (LLMs) can describe and reason about inertial data~\cite{xu2024penetrative,li2025sensorllm,imran2024llasa,yoon2024bymyeyes}, and deployed driver-coaching systems already turn phone telematics into text~\cite{braun2018saferdrive,baig2022drivingbeacon,thompson2026localising}.
Neither line of work is sufficient here.
Sensor--LLM systems interpret activity windows; they do not segment a trip along the road network, and no claim they make is checked against the measurements.
The coaching systems keep content selection in rules and let a language model at most paraphrase~\cite{thompson2026localising}; they do not cover vibration or road surface, and they do not evaluate the factuality of the generated text.
A carefully engineered template narrator is faithful by construction and is a strong baseline~\cite{vandeemter2005templates}, but it can only produce what its designers enumerated.
A trip capsule in our system can take on the order of $10^4$ joint states per road unit before any question, audience or history enters (Section~\ref{sec:related}).

\paragraph{Our position.}
We do not claim that a small language model (SLM) beats a template narrator at summarizing a fixed schema; on that task a template wins on fidelity and cost.
Our thesis is narrower and testable: \emph{a grounded on-device SLM extends a deterministic narrator to outputs that cannot be enumerated at design time---questions users actually ask, audiences no template was written for, evidence states nobody anticipated---and the fidelity it costs can be measured and bounded.}
RoadScribe therefore keeps the template narrator inside the system as a fallback, never lets the language model supply a value, and measures what the model adds over the template.

\paragraph{RoadScribe.}
RoadScribe turns a trip recorded by a phone into a short narrative and answers follow-up questions about it (Figure~\ref{fig:arch}).
Layer~1 estimates the phone's orientation in the vehicle, computes ISO~2631-1-weighted cabin vibration, separates how the ride felt from how the road is, and detects impulsive events with their maneuver context.
Layer~2 matches the trip to OpenStreetMap~\cite{haklay2008osm} and partitions it into \emph{homogeneous road narrative units} (HRNUs) with a change-point detector whose penalty drops at map transitions.
Layer~3 compiles the units into a typed claim ledger and a token-budgeted capsule, and generates text with an on-device SLM under a plan--ground--realize protocol: the model plans over ledger identifiers, deterministic code fills in every value, and a verifier checks each sentence, falling back to the template narrator when it fails.

\paragraph{Contributions.}
The paper makes three novel technical contributions.
\begin{itemize}
\item \textbf{C1 --- Sensing: map-anchored, uncertainty-annotated ride facts from a phone.}
  A Layer~1 pipeline estimates the mount orientation per handling episode with a gated, rotation-equivariant gravity estimate and a closed-form Procrustes heading fit against GNSS; applies the ISO~2631-1 weightings exactly in the frequency domain; reports felt and speed-decoupled road vibration separately with the standard's shock-evaluation rule; and labels impulsive events with a maneuver gate grounded in measured evidence.
  On PVS the heading estimate recovers a rigid mount to within about $2^\circ$ from trip to trip and exposes a $14^\circ$ mounting offset in one vehicle (Section~\ref{sec:l1eval}).
\item \textbf{C2 --- Segmentation: Map-Informed PELT with a corrected pruning rule.}
  We formulate trip segmentation as penalized change-point detection whose penalty is lowered near OpenStreetMap transitions.
  The classical PELT pruning test~\cite{killick2012pelt} assumes a constant penalty; we prove a three-line lemma giving the correct test, $G(t)+C(\mathbf y_{t+1:s})\ge G(s)$, for location-dependent penalties, and extend it to a minimum segment length.
  In 300 random series with map-dependent penalties, the textbook rule returned a sub-optimal partition in 136, while the corrected rule matched exhaustive optimal partitioning in all 300 (Section~\ref{sec:layer2}).
\item \textbf{C3 --- Generation: Plan--Ground--Realize with a verifier and template fallback.}
  The SLM chooses and orders content by emitting a plan constrained by a grammar generated from the capsule's ledger, so it can reference only units and facts that exist; grounding is deterministic; a claim verifier checks every sentence against the ledger, including events, hedges and fields marked missing; and a sentence that fails twice is replaced by the template realization of the same plan item.
  The language model therefore controls selection, ordering, aggregation and phrasing, but not values (Section~\ref{sec:layer3}).
\end{itemize}
In support of these, we contribute a benchmark whose metrics separate language errors from sensing errors (claim precision against the ledger and against ground truth, kinetic--semantic fidelity, anomaly grounding), and the design of a preregistered within-subject study with 28 drivers plus a two-week field deployment to measure usefulness, trust, workload, situation awareness and error detection against strong template baselines (Section~\ref{sec:eval}).
The implementation, with tests that reproduce every number we report as verified, is released as open source.

\paragraph{Results in brief.}
\TBD[: one paragraph with the headline results of E1--E5 and H1--H4 once the benchmark and the study have run. Until then, no result may be stated here.]

The rest of the paper is organized as follows.
Section~\ref{sec:related} positions RoadScribe against sensor--language models, trajectory--language models and data-to-text driver feedback.
Section~\ref{sec:system} formalizes the three layers.
Section~\ref{sec:eval} describes the data, metrics, benchmark and user study.
Section~\ref{sec:discussion} discusses limitations and ethics.
 inside an acmart document (see main.tex).
% =====================================================================
\providecommand{\TBD}[1][]{\textcolor{red}{[TBD#1]}}

\section{Introduction}
\label{sec:intro}

Every smartphone that rides in a car carries an accelerometer, a gyroscope and a GNSS receiver.
Two decades of mobile sensing have used them to monitor roads: early systems detected potholes and braking from phones and dedicated in-vehicle units~\cite{mohan2008nericell,eriksson2008pothole}, and a large body of later work classifies surface types or regresses a roughness index from vehicle vibration~\cite{menegazzo2021pvs}.
The output of these systems is a label or a number per window or per road segment.
The people who drive the roads ask something different.
A commuter wants to know which part of the route was rough and whether to drive it differently.
A passenger prone to car sickness wants to know which stretches and which driving made the ride unpleasant.
A fleet manager wants to know where the vans take shocks, and whether the map is wrong about a surface.
These are questions about a specific trip, in language, and their answers combine facts that a single number cannot hold.

\paragraph{Why a scalar roughness score is the wrong output.}
The International Roughness Index (IRI) is defined on a measured road profile, through a reference quarter-car model driven at a fixed speed~\cite{sayers1986iri}.
A phone in a holder observes something else: the response of an unknown vehicle, through an unknown mount, at whatever speed the driver chose.
Regressing IRI, or any single roughness score, from such data runs into four problems.
\begin{enumerate}
\item \emph{It confounds the road with the drive.} Vibration grows with speed, so a regressor can score well by learning speed. On the PVS dataset~\cite{menegazzo2021pvs}, speed alone separates the three surface types with an AUC of 0.977 \TBD[: re-run audit], and an earlier ``kinetic normalization'' of our own encoded speed almost perfectly (Spearman 0.995 with speed) \TBD[: re-run audit].
\item \emph{It collapses distinct facts.} How the ride \emph{felt} depends on speed; how the \emph{road} is should not. A stretch with one violent speed bump and a stretch of steady cobblestone can share a mean vibration level, yet ISO~2631-1 treats the first as shock-dominated and requires different measures for it~\cite{iso2631}.
\item \emph{It has no place, cause or confidence.} A number does not say where along the street the roughness was, whether a jolt came from the road or from pulling away, or how sure the system is.
\item \emph{It cannot say what is missing.} When the GNSS fix drops out, the gyroscope is absent or the map contradicts the signal, a score is silently wrong; a sentence can say so.
\end{enumerate}

\paragraph{Why not a language model on the raw signal, or a template?}
Large language models (LLMs) can describe and reason about inertial data~\cite{xu2024penetrative,li2025sensorllm,imran2024llasa,yoon2024bymyeyes}, and deployed driver-coaching systems already turn phone telematics into text~\cite{braun2018saferdrive,baig2022drivingbeacon,thompson2026localising}.
Neither line of work is sufficient here.
Sensor--LLM systems interpret activity windows; they do not segment a trip along the road network, and no claim they make is checked against the measurements.
The coaching systems keep content selection in rules and let a language model at most paraphrase~\cite{thompson2026localising}; they do not cover vibration or road surface, and they do not evaluate the factuality of the generated text.
A carefully engineered template narrator is faithful by construction and is a strong baseline~\cite{vandeemter2005templates}, but it can only produce what its designers enumerated.
A trip capsule in our system can take on the order of $10^4$ joint states per road unit before any question, audience or history enters (Section~\ref{sec:related}).

\paragraph{Our position.}
We do not claim that a small language model (SLM) beats a template narrator at summarizing a fixed schema; on that task a template wins on fidelity and cost.
Our thesis is narrower and testable: \emph{a grounded on-device SLM extends a deterministic narrator to outputs that cannot be enumerated at design time---questions users actually ask, audiences no template was written for, evidence states nobody anticipated---and the fidelity it costs can be measured and bounded.}
RoadScribe therefore keeps the template narrator inside the system as a fallback, never lets the language model supply a value, and measures what the model adds over the template.

\paragraph{RoadScribe.}
RoadScribe turns a trip recorded by a phone into a short narrative and answers follow-up questions about it (Figure~\ref{fig:arch}).
Layer~1 estimates the phone's orientation in the vehicle, computes ISO~2631-1-weighted cabin vibration, separates how the ride felt from how the road is, and detects impulsive events with their maneuver context.
Layer~2 matches the trip to OpenStreetMap~\cite{haklay2008osm} and partitions it into \emph{homogeneous road narrative units} (HRNUs) with a change-point detector whose penalty drops at map transitions.
Layer~3 compiles the units into a typed claim ledger and a token-budgeted capsule, and generates text with an on-device SLM under a plan--ground--realize protocol: the model plans over ledger identifiers, deterministic code fills in every value, and a verifier checks each sentence, falling back to the template narrator when it fails.

\paragraph{Contributions.}
The paper makes three novel technical contributions.
\begin{itemize}
\item \textbf{C1 --- Sensing: map-anchored, uncertainty-annotated ride facts from a phone.}
  A Layer~1 pipeline estimates the mount orientation per handling episode with a gated, rotation-equivariant gravity estimate and a closed-form Procrustes heading fit against GNSS; applies the ISO~2631-1 weightings exactly in the frequency domain; reports felt and speed-decoupled road vibration separately with the standard's shock-evaluation rule; and labels impulsive events with a maneuver gate grounded in measured evidence.
  On PVS the heading estimate recovers a rigid mount to within about $2^\circ$ from trip to trip and exposes a $14^\circ$ mounting offset in one vehicle (Section~\ref{sec:l1eval}).
\item \textbf{C2 --- Segmentation: Map-Informed PELT with a corrected pruning rule.}
  We formulate trip segmentation as penalized change-point detection whose penalty is lowered near OpenStreetMap transitions.
  The classical PELT pruning test~\cite{killick2012pelt} assumes a constant penalty; we prove a three-line lemma giving the correct test, $G(t)+C(\mathbf y_{t+1:s})\ge G(s)$, for location-dependent penalties, and extend it to a minimum segment length.
  In 300 random series with map-dependent penalties, the textbook rule returned a sub-optimal partition in 136, while the corrected rule matched exhaustive optimal partitioning in all 300 (Section~\ref{sec:layer2}).
\item \textbf{C3 --- Generation: Plan--Ground--Realize with a verifier and template fallback.}
  The SLM chooses and orders content by emitting a plan constrained by a grammar generated from the capsule's ledger, so it can reference only units and facts that exist; grounding is deterministic; a claim verifier checks every sentence against the ledger, including events, hedges and fields marked missing; and a sentence that fails twice is replaced by the template realization of the same plan item.
  The language model therefore controls selection, ordering, aggregation and phrasing, but not values (Section~\ref{sec:layer3}).
\end{itemize}
In support of these, we contribute a benchmark whose metrics separate language errors from sensing errors (claim precision against the ledger and against ground truth, kinetic--semantic fidelity, anomaly grounding), and the design of a preregistered within-subject study with 28 drivers plus a two-week field deployment to measure usefulness, trust, workload, situation awareness and error detection against strong template baselines (Section~\ref{sec:eval}).
The implementation, with tests that reproduce every number we report as verified, is released as open source.

\paragraph{Results in brief.}
\TBD[: one paragraph with the headline results of E1--E5 and H1--H4 once the benchmark and the study have run. Until then, no result may be stated here.]

The rest of the paper is organized as follows.
Section~\ref{sec:related} positions RoadScribe against sensor--language models, trajectory--language models and data-to-text driver feedback.
Section~\ref{sec:system} formalizes the three layers.
Section~\ref{sec:eval} describes the data, metrics, benchmark and user study.
Section~\ref{sec:discussion} discusses limitations and ethics.
 inside an acmart document (see main.tex).
% =====================================================================
\providecommand{\TBD}[1][]{\textcolor{red}{[TBD#1]}}

\section{Introduction}
\label{sec:intro}

Every smartphone that rides in a car carries an accelerometer, a gyroscope and a GNSS receiver.
Two decades of mobile sensing have used them to monitor roads: early systems detected potholes and braking from phones and dedicated in-vehicle units~\cite{mohan2008nericell,eriksson2008pothole}, and a large body of later work classifies surface types or regresses a roughness index from vehicle vibration~\cite{menegazzo2021pvs}.
The output of these systems is a label or a number per window or per road segment.
The people who drive the roads ask something different.
A commuter wants to know which part of the route was rough and whether to drive it differently.
A passenger prone to car sickness wants to know which stretches and which driving made the ride unpleasant.
A fleet manager wants to know where the vans take shocks, and whether the map is wrong about a surface.
These are questions about a specific trip, in language, and their answers combine facts that a single number cannot hold.

\paragraph{Why a scalar roughness score is the wrong output.}
The International Roughness Index (IRI) is defined on a measured road profile, through a reference quarter-car model driven at a fixed speed~\cite{sayers1986iri}.
A phone in a holder observes something else: the response of an unknown vehicle, through an unknown mount, at whatever speed the driver chose.
Regressing IRI, or any single roughness score, from such data runs into four problems.
\begin{enumerate}
\item \emph{It confounds the road with the drive.} Vibration grows with speed, so a regressor can score well by learning speed. On the PVS dataset~\cite{menegazzo2021pvs}, speed alone separates the three surface types with an AUC of 0.977 \TBD[: re-run audit], and an earlier ``kinetic normalization'' of our own encoded speed almost perfectly (Spearman 0.995 with speed) \TBD[: re-run audit].
\item \emph{It collapses distinct facts.} How the ride \emph{felt} depends on speed; how the \emph{road} is should not. A stretch with one violent speed bump and a stretch of steady cobblestone can share a mean vibration level, yet ISO~2631-1 treats the first as shock-dominated and requires different measures for it~\cite{iso2631}.
\item \emph{It has no place, cause or confidence.} A number does not say where along the street the roughness was, whether a jolt came from the road or from pulling away, or how sure the system is.
\item \emph{It cannot say what is missing.} When the GNSS fix drops out, the gyroscope is absent or the map contradicts the signal, a score is silently wrong; a sentence can say so.
\end{enumerate}

\paragraph{Why not a language model on the raw signal, or a template?}
Large language models (LLMs) can describe and reason about inertial data~\cite{xu2024penetrative,li2025sensorllm,imran2024llasa,yoon2024bymyeyes}, and deployed driver-coaching systems already turn phone telematics into text~\cite{braun2018saferdrive,baig2022drivingbeacon,thompson2026localising}.
Neither line of work is sufficient here.
Sensor--LLM systems interpret activity windows; they do not segment a trip along the road network, and no claim they make is checked against the measurements.
The coaching systems keep content selection in rules and let a language model at most paraphrase~\cite{thompson2026localising}; they do not cover vibration or road surface, and they do not evaluate the factuality of the generated text.
A carefully engineered template narrator is faithful by construction and is a strong baseline~\cite{vandeemter2005templates}, but it can only produce what its designers enumerated.
A trip capsule in our system can take on the order of $10^4$ joint states per road unit before any question, audience or history enters (Section~\ref{sec:related}).

\paragraph{Our position.}
We do not claim that a small language model (SLM) beats a template narrator at summarizing a fixed schema; on that task a template wins on fidelity and cost.
Our thesis is narrower and testable: \emph{a grounded on-device SLM extends a deterministic narrator to outputs that cannot be enumerated at design time---questions users actually ask, audiences no template was written for, evidence states nobody anticipated---and the fidelity it costs can be measured and bounded.}
RoadScribe therefore keeps the template narrator inside the system as a fallback, never lets the language model supply a value, and measures what the model adds over the template.

\paragraph{RoadScribe.}
RoadScribe turns a trip recorded by a phone into a short narrative and answers follow-up questions about it (Figure~\ref{fig:arch}).
Layer~1 estimates the phone's orientation in the vehicle, computes ISO~2631-1-weighted cabin vibration, separates how the ride felt from how the road is, and detects impulsive events with their maneuver context.
Layer~2 matches the trip to OpenStreetMap~\cite{haklay2008osm} and partitions it into \emph{homogeneous road narrative units} (HRNUs) with a change-point detector whose penalty drops at map transitions.
Layer~3 compiles the units into a typed claim ledger and a token-budgeted capsule, and generates text with an on-device SLM under a plan--ground--realize protocol: the model plans over ledger identifiers, deterministic code fills in every value, and a verifier checks each sentence, falling back to the template narrator when it fails.

\paragraph{Contributions.}
The paper makes three novel technical contributions.
\begin{itemize}
\item \textbf{C1 --- Sensing: map-anchored, uncertainty-annotated ride facts from a phone.}
  A Layer~1 pipeline estimates the mount orientation per handling episode with a gated, rotation-equivariant gravity estimate and a closed-form Procrustes heading fit against GNSS; applies the ISO~2631-1 weightings exactly in the frequency domain; reports felt and speed-decoupled road vibration separately with the standard's shock-evaluation rule; and labels impulsive events with a maneuver gate grounded in measured evidence.
  On PVS the heading estimate recovers a rigid mount to within about $2^\circ$ from trip to trip and exposes a $14^\circ$ mounting offset in one vehicle (Section~\ref{sec:l1eval}).
\item \textbf{C2 --- Segmentation: Map-Informed PELT with a corrected pruning rule.}
  We formulate trip segmentation as penalized change-point detection whose penalty is lowered near OpenStreetMap transitions.
  The classical PELT pruning test~\cite{killick2012pelt} assumes a constant penalty; we prove a three-line lemma giving the correct test, $G(t)+C(\mathbf y_{t+1:s})\ge G(s)$, for location-dependent penalties, and extend it to a minimum segment length.
  In 300 random series with map-dependent penalties, the textbook rule returned a sub-optimal partition in 136, while the corrected rule matched exhaustive optimal partitioning in all 300 (Section~\ref{sec:layer2}).
\item \textbf{C3 --- Generation: Plan--Ground--Realize with a verifier and template fallback.}
  The SLM chooses and orders content by emitting a plan constrained by a grammar generated from the capsule's ledger, so it can reference only units and facts that exist; grounding is deterministic; a claim verifier checks every sentence against the ledger, including events, hedges and fields marked missing; and a sentence that fails twice is replaced by the template realization of the same plan item.
  The language model therefore controls selection, ordering, aggregation and phrasing, but not values (Section~\ref{sec:layer3}).
\end{itemize}
In support of these, we contribute a benchmark whose metrics separate language errors from sensing errors (claim precision against the ledger and against ground truth, kinetic--semantic fidelity, anomaly grounding), and the design of a preregistered within-subject study with 28 drivers plus a two-week field deployment to measure usefulness, trust, workload, situation awareness and error detection against strong template baselines (Section~\ref{sec:eval}).
The implementation, with tests that reproduce every number we report as verified, is released as open source.

\paragraph{Results in brief.}
\TBD[: one paragraph with the headline results of E1--E5 and H1--H4 once the benchmark and the study have run. Until then, no result may be stated here.]

The rest of the paper is organized as follows.
Section~\ref{sec:related} positions RoadScribe against sensor--language models, trajectory--language models and data-to-text driver feedback.
Section~\ref{sec:system} formalizes the three layers.
Section~\ref{sec:eval} describes the data, metrics, benchmark and user study.
Section~\ref{sec:discussion} discusses limitations and ethics.

% =====================================================================
% Section 2 -- Related Work
% Usage: % =====================================================================
% Section 2 -- Related Work
% Usage: % =====================================================================
% Section 2 -- Related Work
% Usage: \input{sec2_related} inside an acmart document (see main.tex).
% =====================================================================

\section{Related Work}
\label{sec:related}

RoadScribe sits where three lines of work meet without overlapping: language models that read sensor streams (Stream~A), language models that read trajectories (Stream~B), and data-to-text systems that give drivers feedback (Stream~C).
We review each, then road sensing and template generation, and close with four limitations that no prior system addresses together.

\subsection{Stream A: Language Models over Sensor Data}

\emph{Penetrative AI} showed that a general LLM, prompted with raw or lightly processed IoT readings, can carry out perception tasks such as activity recognition and heartbeat detection using its world knowledge~\cite{xu2024penetrative}.
\emph{SensorLLM} aligns an LLM with motion-sensor time series in two stages: it first learns to describe trends in the channels from automatically generated text, then is tuned for human activity recognition, matching or exceeding dedicated models on several datasets~\cite{li2025sensorllm}.
\emph{LLaSA} couples an IMU encoder with an LLM and releases narrations and question-answer data for reasoning about activities~\cite{imran2024llasa}, and \emph{By My Eyes} renders sensor streams as plots for a multimodal LLM, improving accuracy and cutting token cost relative to textual prompts~\cite{yoon2024bymyeyes}.

These systems establish that language models can interpret inertial data.
Three properties separate them from our setting.
Their unit of analysis is an activity window; none partitions a long recording along a spatial network or anchors statements to places.
Their outputs are labels or free text whose numbers and events are not checked against the measurements at the level of individual claims; evaluation uses classification accuracy, text similarity or judged plausibility.
And the language model reads the signal directly, so any value it states may have been invented.
RoadScribe inverts the last point: signal processing produces every value, and the language model only selects and phrases ledger facts.

\subsection{Stream B: Language Models over Trajectories}

A second line connects language models to movement data.
\emph{Traj-LLM} adapts a pre-trained LLM, with parameter-efficient fine-tuning and a lane-aware module, to predict the future trajectories of traffic agents~\cite{lan2024trajllm}; here the language model is a sequence learner, and its output is a trajectory, not text.
Benchmarks such as \emph{TrajPrism} pair urban GPS trajectories with language for instruction-conditioned generation, retrieval and trajectory captioning, and report a large gap between geometry-only models and language-grounded tasks~\cite{li2026trajprism}.
Earlier, symbolic generators produced route descriptions from map paths~\cite{dale2005route}.

Trajectory--language work describes \emph{where} a vehicle went.
It does not observe \emph{what the road did to the vehicle}: GNSS traces carry no vibration, so these systems cannot say that a street was rough, that a jolt came from a speed bump or from pulling away, or that the map's surface tag is wrong.
RoadScribe uses the trajectory only as a spatial scaffold, through map matching, and takes its content from inertial sensing.

\subsection{Stream C: Data-to-Text Feedback for Drivers}

Data-to-text generation has a long record of turning sensor data into prose, from intensive-care summaries~\cite{portet2009babytalk} to weather forecasts~\cite{reiter2005sumtime}; the classic pipeline separates signal analysis, data interpretation, document planning and realization~\cite{gatt2018survey}.
For drivers, \emph{SaferDrive} turned insurance telematics into weekly rule-based feedback on speeding and harsh events~\cite{braun2018saferdrive}.
\emph{DrivingBeacon} extended this with phone use and geographic context and evaluated it in a four-week A/B test with eight drivers~\cite{baig2022drivingbeacon}.
Most recently, Thompson et al. compared two localized systems~\cite{thompson2026localising}: a UK system for post-trip reflection with road and place context, deployed with 30 drivers for 12 weeks, and a Nigerian system for safety education.
In both, structured components decide what may be said and a cloud LLM may only adapt how it is said; the authors state that they do not independently evaluate the factuality of the generated text, and there is no controlled comparison between the template and the LLM variants.

This stream is the closest to RoadScribe, and we adopt its central principle: content selection must be controlled.
It differs from our setting in content (driving behaviour rather than road and ride), in where the model runs (cloud rather than device), and in evaluation (engagement and behaviour change rather than claim-level fidelity).

\subsection{Road Sensing and Template Generation}

Nericell~\cite{mohan2008nericell} and Pothole Patrol~\cite{eriksson2008pothole} established phone and in-vehicle accelerometer road monitoring, including virtual reorientation of a phone to the vehicle frame.
Many later studies classify surface type or regress roughness from inertial data~\cite{menegazzo2021pvs}.
Their outputs are labels or indices.
We reuse their sensing insights but treat measurement uncertainty, maneuver confounds and speed explicitly (Section~\ref{sec:layer1}), and we evaluate against speed-only baselines because surface type and speed are confounded in driving data.

On the language side, van Deemter et al. argued that the opposition between template and ``real'' generation is false: a template system can be linguistically rich~\cite{vandeemter2005templates}.
We take that argument seriously.
Our template baseline is a full pipeline with content selection, aggregation, referring expressions and a fallback for every missing field, and it is also RoadScribe's fallback.
The question is where templates stop scaling.
Per HRNU, a capsule combines a felt class (6 values), a road class (6), a shock flag (2), an event pattern (about 4), map--signal agreement (2), a confidence band (3) and a missing-channel pattern ($2^3$): about 13{,}800 joint states before an audience, a question or a trip history is added.
Most need no special wording; the ones that do are interactions, such as a stretch that felt very rough mostly because of speed on a surface that rides like cobblestone but is mapped as asphalt.
Our benchmark measures this coverage gap instead of asserting it (Section~\ref{sec:bench}).

Faithfulness is a known weakness of neural generation~\cite{ji2023hallucination}.
We evaluate it with atomic claims~\cite{min2023factscore,wiseman2017challenges}, express uncertainty with calibrated hedges~\cite{kent1964wep,mielke2022linguistic}, and constrain the plan with grammar-based decoding~\cite{geng2023grammar,ggerganov_llamacpp}, which removes the output-format failures reported for on-device SLMs in shipped apps~\cite{oliveira2026lessismore}.

\subsection{Four Limitations of the State of the Art}
\label{sec:limitations}

Taken together, prior work leaves four gaps that RoadScribe is designed to close.

\paragraph{L1: No claim-level grounding.}
Sensor--language models state values they read or infer from the signal, and the coaching systems do not evaluate the factuality of their text~\cite{thompson2026localising}.
No system checks each number, event, place and severity word against the measurements.
\emph{RoadScribe} keeps every value in a typed ledger, lets the model reference only ledger identifiers, and verifies each sentence against the ledger before release (Section~\ref{sec:layer3}).
Its metrics then separate language errors from sensing errors by scoring claims against both the ledger and ground truth (Section~\ref{sec:metrics}).

\paragraph{L2: No road-anchored segmentation of ride content.}
Sensor--language models reason over activity windows, trajectory--language models over paths without vibration, and road-monitoring systems over fixed windows or map segments.
None partitions a trip into units that are homogeneous in ride content \emph{and} aligned with the road network.
\emph{RoadScribe} segments in the distance domain with a penalty that drops at map transitions, with a pruning rule that is provably exact for such penalties (Section~\ref{sec:layer2}).

\paragraph{L3: Uncertainty and missing evidence never reach the text.}
Existing systems report what they detect but not how sure they are or what they could not observe.
A jolt during stop-and-go traffic, a 48-second GNSS gap or a map tag that contradicts the signal all disappear into fluent text.
\emph{RoadScribe} carries event confidence, map-match confidence and \textsc{missing} facts with reasons into the ledger, maps confidence to hedges by a fixed policy, and counts any value stated for a missing field as a hallucination (Sections~\ref{sec:layer3} and~\ref{sec:metrics}).

\paragraph{L4: No controlled comparison with templates, and no on-device evaluation.}
The coaching systems compare neither their LLM variants with their template variants nor run the model on the phone; sensor--language work uses large models off-device.
Whether a language model adds anything a template cannot, and at what fidelity, latency and energy cost on a phone, is unmeasured.
\emph{RoadScribe} is designed to run a quantized SLM on device, keeps the template narrator as both baseline and fallback, and specifies six narrators, including template-plus-variation and template-plus-paraphrase controls, on tasks designed to expose where templates run out (Section~\ref{sec:eval}).
 inside an acmart document (see main.tex).
% =====================================================================

\section{Related Work}
\label{sec:related}

RoadScribe sits where three lines of work meet without overlapping: language models that read sensor streams (Stream~A), language models that read trajectories (Stream~B), and data-to-text systems that give drivers feedback (Stream~C).
We review each, then road sensing and template generation, and close with four limitations that no prior system addresses together.

\subsection{Stream A: Language Models over Sensor Data}

\emph{Penetrative AI} showed that a general LLM, prompted with raw or lightly processed IoT readings, can carry out perception tasks such as activity recognition and heartbeat detection using its world knowledge~\cite{xu2024penetrative}.
\emph{SensorLLM} aligns an LLM with motion-sensor time series in two stages: it first learns to describe trends in the channels from automatically generated text, then is tuned for human activity recognition, matching or exceeding dedicated models on several datasets~\cite{li2025sensorllm}.
\emph{LLaSA} couples an IMU encoder with an LLM and releases narrations and question-answer data for reasoning about activities~\cite{imran2024llasa}, and \emph{By My Eyes} renders sensor streams as plots for a multimodal LLM, improving accuracy and cutting token cost relative to textual prompts~\cite{yoon2024bymyeyes}.

These systems establish that language models can interpret inertial data.
Three properties separate them from our setting.
Their unit of analysis is an activity window; none partitions a long recording along a spatial network or anchors statements to places.
Their outputs are labels or free text whose numbers and events are not checked against the measurements at the level of individual claims; evaluation uses classification accuracy, text similarity or judged plausibility.
And the language model reads the signal directly, so any value it states may have been invented.
RoadScribe inverts the last point: signal processing produces every value, and the language model only selects and phrases ledger facts.

\subsection{Stream B: Language Models over Trajectories}

A second line connects language models to movement data.
\emph{Traj-LLM} adapts a pre-trained LLM, with parameter-efficient fine-tuning and a lane-aware module, to predict the future trajectories of traffic agents~\cite{lan2024trajllm}; here the language model is a sequence learner, and its output is a trajectory, not text.
Benchmarks such as \emph{TrajPrism} pair urban GPS trajectories with language for instruction-conditioned generation, retrieval and trajectory captioning, and report a large gap between geometry-only models and language-grounded tasks~\cite{li2026trajprism}.
Earlier, symbolic generators produced route descriptions from map paths~\cite{dale2005route}.

Trajectory--language work describes \emph{where} a vehicle went.
It does not observe \emph{what the road did to the vehicle}: GNSS traces carry no vibration, so these systems cannot say that a street was rough, that a jolt came from a speed bump or from pulling away, or that the map's surface tag is wrong.
RoadScribe uses the trajectory only as a spatial scaffold, through map matching, and takes its content from inertial sensing.

\subsection{Stream C: Data-to-Text Feedback for Drivers}

Data-to-text generation has a long record of turning sensor data into prose, from intensive-care summaries~\cite{portet2009babytalk} to weather forecasts~\cite{reiter2005sumtime}; the classic pipeline separates signal analysis, data interpretation, document planning and realization~\cite{gatt2018survey}.
For drivers, \emph{SaferDrive} turned insurance telematics into weekly rule-based feedback on speeding and harsh events~\cite{braun2018saferdrive}.
\emph{DrivingBeacon} extended this with phone use and geographic context and evaluated it in a four-week A/B test with eight drivers~\cite{baig2022drivingbeacon}.
Most recently, Thompson et al. compared two localized systems~\cite{thompson2026localising}: a UK system for post-trip reflection with road and place context, deployed with 30 drivers for 12 weeks, and a Nigerian system for safety education.
In both, structured components decide what may be said and a cloud LLM may only adapt how it is said; the authors state that they do not independently evaluate the factuality of the generated text, and there is no controlled comparison between the template and the LLM variants.

This stream is the closest to RoadScribe, and we adopt its central principle: content selection must be controlled.
It differs from our setting in content (driving behaviour rather than road and ride), in where the model runs (cloud rather than device), and in evaluation (engagement and behaviour change rather than claim-level fidelity).

\subsection{Road Sensing and Template Generation}

Nericell~\cite{mohan2008nericell} and Pothole Patrol~\cite{eriksson2008pothole} established phone and in-vehicle accelerometer road monitoring, including virtual reorientation of a phone to the vehicle frame.
Many later studies classify surface type or regress roughness from inertial data~\cite{menegazzo2021pvs}.
Their outputs are labels or indices.
We reuse their sensing insights but treat measurement uncertainty, maneuver confounds and speed explicitly (Section~\ref{sec:layer1}), and we evaluate against speed-only baselines because surface type and speed are confounded in driving data.

On the language side, van Deemter et al. argued that the opposition between template and ``real'' generation is false: a template system can be linguistically rich~\cite{vandeemter2005templates}.
We take that argument seriously.
Our template baseline is a full pipeline with content selection, aggregation, referring expressions and a fallback for every missing field, and it is also RoadScribe's fallback.
The question is where templates stop scaling.
Per HRNU, a capsule combines a felt class (6 values), a road class (6), a shock flag (2), an event pattern (about 4), map--signal agreement (2), a confidence band (3) and a missing-channel pattern ($2^3$): about 13{,}800 joint states before an audience, a question or a trip history is added.
Most need no special wording; the ones that do are interactions, such as a stretch that felt very rough mostly because of speed on a surface that rides like cobblestone but is mapped as asphalt.
Our benchmark measures this coverage gap instead of asserting it (Section~\ref{sec:bench}).

Faithfulness is a known weakness of neural generation~\cite{ji2023hallucination}.
We evaluate it with atomic claims~\cite{min2023factscore,wiseman2017challenges}, express uncertainty with calibrated hedges~\cite{kent1964wep,mielke2022linguistic}, and constrain the plan with grammar-based decoding~\cite{geng2023grammar,ggerganov_llamacpp}, which removes the output-format failures reported for on-device SLMs in shipped apps~\cite{oliveira2026lessismore}.

\subsection{Four Limitations of the State of the Art}
\label{sec:limitations}

Taken together, prior work leaves four gaps that RoadScribe is designed to close.

\paragraph{L1: No claim-level grounding.}
Sensor--language models state values they read or infer from the signal, and the coaching systems do not evaluate the factuality of their text~\cite{thompson2026localising}.
No system checks each number, event, place and severity word against the measurements.
\emph{RoadScribe} keeps every value in a typed ledger, lets the model reference only ledger identifiers, and verifies each sentence against the ledger before release (Section~\ref{sec:layer3}).
Its metrics then separate language errors from sensing errors by scoring claims against both the ledger and ground truth (Section~\ref{sec:metrics}).

\paragraph{L2: No road-anchored segmentation of ride content.}
Sensor--language models reason over activity windows, trajectory--language models over paths without vibration, and road-monitoring systems over fixed windows or map segments.
None partitions a trip into units that are homogeneous in ride content \emph{and} aligned with the road network.
\emph{RoadScribe} segments in the distance domain with a penalty that drops at map transitions, with a pruning rule that is provably exact for such penalties (Section~\ref{sec:layer2}).

\paragraph{L3: Uncertainty and missing evidence never reach the text.}
Existing systems report what they detect but not how sure they are or what they could not observe.
A jolt during stop-and-go traffic, a 48-second GNSS gap or a map tag that contradicts the signal all disappear into fluent text.
\emph{RoadScribe} carries event confidence, map-match confidence and \textsc{missing} facts with reasons into the ledger, maps confidence to hedges by a fixed policy, and counts any value stated for a missing field as a hallucination (Sections~\ref{sec:layer3} and~\ref{sec:metrics}).

\paragraph{L4: No controlled comparison with templates, and no on-device evaluation.}
The coaching systems compare neither their LLM variants with their template variants nor run the model on the phone; sensor--language work uses large models off-device.
Whether a language model adds anything a template cannot, and at what fidelity, latency and energy cost on a phone, is unmeasured.
\emph{RoadScribe} is designed to run a quantized SLM on device, keeps the template narrator as both baseline and fallback, and specifies six narrators, including template-plus-variation and template-plus-paraphrase controls, on tasks designed to expose where templates run out (Section~\ref{sec:eval}).
 inside an acmart document (see main.tex).
% =====================================================================

\section{Related Work}
\label{sec:related}

RoadScribe sits where three lines of work meet without overlapping: language models that read sensor streams (Stream~A), language models that read trajectories (Stream~B), and data-to-text systems that give drivers feedback (Stream~C).
We review each, then road sensing and template generation, and close with four limitations that no prior system addresses together.

\subsection{Stream A: Language Models over Sensor Data}

\emph{Penetrative AI} showed that a general LLM, prompted with raw or lightly processed IoT readings, can carry out perception tasks such as activity recognition and heartbeat detection using its world knowledge~\cite{xu2024penetrative}.
\emph{SensorLLM} aligns an LLM with motion-sensor time series in two stages: it first learns to describe trends in the channels from automatically generated text, then is tuned for human activity recognition, matching or exceeding dedicated models on several datasets~\cite{li2025sensorllm}.
\emph{LLaSA} couples an IMU encoder with an LLM and releases narrations and question-answer data for reasoning about activities~\cite{imran2024llasa}, and \emph{By My Eyes} renders sensor streams as plots for a multimodal LLM, improving accuracy and cutting token cost relative to textual prompts~\cite{yoon2024bymyeyes}.

These systems establish that language models can interpret inertial data.
Three properties separate them from our setting.
Their unit of analysis is an activity window; none partitions a long recording along a spatial network or anchors statements to places.
Their outputs are labels or free text whose numbers and events are not checked against the measurements at the level of individual claims; evaluation uses classification accuracy, text similarity or judged plausibility.
And the language model reads the signal directly, so any value it states may have been invented.
RoadScribe inverts the last point: signal processing produces every value, and the language model only selects and phrases ledger facts.

\subsection{Stream B: Language Models over Trajectories}

A second line connects language models to movement data.
\emph{Traj-LLM} adapts a pre-trained LLM, with parameter-efficient fine-tuning and a lane-aware module, to predict the future trajectories of traffic agents~\cite{lan2024trajllm}; here the language model is a sequence learner, and its output is a trajectory, not text.
Benchmarks such as \emph{TrajPrism} pair urban GPS trajectories with language for instruction-conditioned generation, retrieval and trajectory captioning, and report a large gap between geometry-only models and language-grounded tasks~\cite{li2026trajprism}.
Earlier, symbolic generators produced route descriptions from map paths~\cite{dale2005route}.

Trajectory--language work describes \emph{where} a vehicle went.
It does not observe \emph{what the road did to the vehicle}: GNSS traces carry no vibration, so these systems cannot say that a street was rough, that a jolt came from a speed bump or from pulling away, or that the map's surface tag is wrong.
RoadScribe uses the trajectory only as a spatial scaffold, through map matching, and takes its content from inertial sensing.

\subsection{Stream C: Data-to-Text Feedback for Drivers}

Data-to-text generation has a long record of turning sensor data into prose, from intensive-care summaries~\cite{portet2009babytalk} to weather forecasts~\cite{reiter2005sumtime}; the classic pipeline separates signal analysis, data interpretation, document planning and realization~\cite{gatt2018survey}.
For drivers, \emph{SaferDrive} turned insurance telematics into weekly rule-based feedback on speeding and harsh events~\cite{braun2018saferdrive}.
\emph{DrivingBeacon} extended this with phone use and geographic context and evaluated it in a four-week A/B test with eight drivers~\cite{baig2022drivingbeacon}.
Most recently, Thompson et al. compared two localized systems~\cite{thompson2026localising}: a UK system for post-trip reflection with road and place context, deployed with 30 drivers for 12 weeks, and a Nigerian system for safety education.
In both, structured components decide what may be said and a cloud LLM may only adapt how it is said; the authors state that they do not independently evaluate the factuality of the generated text, and there is no controlled comparison between the template and the LLM variants.

This stream is the closest to RoadScribe, and we adopt its central principle: content selection must be controlled.
It differs from our setting in content (driving behaviour rather than road and ride), in where the model runs (cloud rather than device), and in evaluation (engagement and behaviour change rather than claim-level fidelity).

\subsection{Road Sensing and Template Generation}

Nericell~\cite{mohan2008nericell} and Pothole Patrol~\cite{eriksson2008pothole} established phone and in-vehicle accelerometer road monitoring, including virtual reorientation of a phone to the vehicle frame.
Many later studies classify surface type or regress roughness from inertial data~\cite{menegazzo2021pvs}.
Their outputs are labels or indices.
We reuse their sensing insights but treat measurement uncertainty, maneuver confounds and speed explicitly (Section~\ref{sec:layer1}), and we evaluate against speed-only baselines because surface type and speed are confounded in driving data.

On the language side, van Deemter et al. argued that the opposition between template and ``real'' generation is false: a template system can be linguistically rich~\cite{vandeemter2005templates}.
We take that argument seriously.
Our template baseline is a full pipeline with content selection, aggregation, referring expressions and a fallback for every missing field, and it is also RoadScribe's fallback.
The question is where templates stop scaling.
Per HRNU, a capsule combines a felt class (6 values), a road class (6), a shock flag (2), an event pattern (about 4), map--signal agreement (2), a confidence band (3) and a missing-channel pattern ($2^3$): about 13{,}800 joint states before an audience, a question or a trip history is added.
Most need no special wording; the ones that do are interactions, such as a stretch that felt very rough mostly because of speed on a surface that rides like cobblestone but is mapped as asphalt.
Our benchmark measures this coverage gap instead of asserting it (Section~\ref{sec:bench}).

Faithfulness is a known weakness of neural generation~\cite{ji2023hallucination}.
We evaluate it with atomic claims~\cite{min2023factscore,wiseman2017challenges}, express uncertainty with calibrated hedges~\cite{kent1964wep,mielke2022linguistic}, and constrain the plan with grammar-based decoding~\cite{geng2023grammar,ggerganov_llamacpp}, which removes the output-format failures reported for on-device SLMs in shipped apps~\cite{oliveira2026lessismore}.

\subsection{Four Limitations of the State of the Art}
\label{sec:limitations}

Taken together, prior work leaves four gaps that RoadScribe is designed to close.

\paragraph{L1: No claim-level grounding.}
Sensor--language models state values they read or infer from the signal, and the coaching systems do not evaluate the factuality of their text~\cite{thompson2026localising}.
No system checks each number, event, place and severity word against the measurements.
\emph{RoadScribe} keeps every value in a typed ledger, lets the model reference only ledger identifiers, and verifies each sentence against the ledger before release (Section~\ref{sec:layer3}).
Its metrics then separate language errors from sensing errors by scoring claims against both the ledger and ground truth (Section~\ref{sec:metrics}).

\paragraph{L2: No road-anchored segmentation of ride content.}
Sensor--language models reason over activity windows, trajectory--language models over paths without vibration, and road-monitoring systems over fixed windows or map segments.
None partitions a trip into units that are homogeneous in ride content \emph{and} aligned with the road network.
\emph{RoadScribe} segments in the distance domain with a penalty that drops at map transitions, with a pruning rule that is provably exact for such penalties (Section~\ref{sec:layer2}).

\paragraph{L3: Uncertainty and missing evidence never reach the text.}
Existing systems report what they detect but not how sure they are or what they could not observe.
A jolt during stop-and-go traffic, a 48-second GNSS gap or a map tag that contradicts the signal all disappear into fluent text.
\emph{RoadScribe} carries event confidence, map-match confidence and \textsc{missing} facts with reasons into the ledger, maps confidence to hedges by a fixed policy, and counts any value stated for a missing field as a hallucination (Sections~\ref{sec:layer3} and~\ref{sec:metrics}).

\paragraph{L4: No controlled comparison with templates, and no on-device evaluation.}
The coaching systems compare neither their LLM variants with their template variants nor run the model on the phone; sensor--language work uses large models off-device.
Whether a language model adds anything a template cannot, and at what fidelity, latency and energy cost on a phone, is unmeasured.
\emph{RoadScribe} is designed to run a quantized SLM on device, keeps the template narrator as both baseline and fallback, and specifies six narrators, including template-plus-variation and template-plus-paraphrase controls, on tasks designed to expose where templates run out (Section~\ref{sec:eval}).

% =====================================================================
% Section 3 -- The RoadScribe system
% Usage: % =====================================================================
% Section 3 -- The RoadScribe system
% Usage: % =====================================================================
% Section 3 -- The RoadScribe system
% Usage: \input{sec3_system} inside an acmart document.
% Preamble requirements (acmart already loads amsmath, amsthm, booktabs, xcolor):
%   \usepackage{algorithm}
%   \usepackage{algpseudocode}
% Bibliography keys used below (add to your .bib):
%   iso2631, iso8855, iso8608, iso13473, newson2009hmm, killick2012pelt,
%   haynes2017crops, kaiser1990tkeo, mohan2008nericell, eriksson2008pothole,
%   menegazzo2021pvs, ggerganov_llamacpp
% =====================================================================
\providecommand{\R}{\mathbb{R}}
\providecommand{\ind}{\mathbf{1}}
\providecommand{\TBD}[1][]{\textcolor{red}{[TBD#1]}}

\section{The RoadScribe System}
\label{sec:system}

RoadScribe turns a trip recorded by a commodity smartphone into a short narrative and answers follow-up questions about it.
The pipeline has three layers (Figure~\ref{fig:arch}).
Layer~1 converts inertial and GNSS samples into vehicle-frame ride measures and candidate events.
Layer~2 anchors these measures to the road network and partitions the trip into homogeneous road narrative units (HRNUs).
Layer~3 compiles the units into a token-budgeted capsule and a machine-checkable claim ledger, and generates text with an on-device small language model (SLM) under a plan--ground--realize protocol.

Four principles shape the design.
(P1)~Every number, place name and event in the output can be traced to a ledger fact.
(P2)~The language model decides what to say and how to say it, but never supplies a value; values come from the ledger.
(P3)~Uncertainty estimated in Layers~1--2 reaches the text as calibrated hedges.
(P4)~When verification fails, a deterministic realizer takes over, so the SLM can add coverage but cannot remove faithfulness.

\begin{figure}[t]
  \centering
  \fbox{\parbox{0.95\linewidth}{\centering\vspace{2em}Architecture figure placeholder: IMU+GNSS $\rightarrow$ Layer 1 (orientation, ISO weighting, events) $\rightarrow$ Layer 2 (map matching, MI-PELT, HRNUs) $\rightarrow$ Layer 3 (ledger, KST capsule, plan--ground--realize, verifier, template fallback) $\rightarrow$ narrative / answers\vspace{2em}}}
  \caption{RoadScribe architecture.}
  \label{fig:arch}
\end{figure}

% ---------------------------------------------------------------------
\subsection{Layer 1: Kinetic Signal Processing}
\label{sec:layer1}

\paragraph{Notation.}
The IMU delivers specific force $\mathbf f^b[n]\in\R^3$ and angular rate $\boldsymbol\omega^b[n]\in\R^3$ in the phone body frame $b$ at rate $f_s$ ($T_s=1/f_s$).
GNSS delivers speed $v_k$ at times $t_k$ (1~Hz).
The vehicle frame $v$ follows ISO~8855~\cite{iso8855} ($x$ forward, $y$ left, $z$ up).
We estimate $\mathbf R_{vb}\in SO(3)$ and the gravity-free acceleration
\begin{equation}
  \mathbf a^v[n]=\mathbf R_{vb}\,\mathbf f^b[n]-g\,\mathbf e_z .
  \label{eq:av}
\end{equation}
Speed is interpolated to IMU time with a monotone cubic (PCHIP), giving $v(t)$ and $a^G_{\mathrm{lon}}(t)=\dot v(t)$.

\paragraph{Mount episodes.}
A phone in a holder is static relative to the vehicle except when handled.
We mark candidate handling where the gyroscope reports rotation that the vehicle does not perform, $\lVert\bar{\boldsymbol\omega}^b[n]\rVert-|\dot\psi^G(t_n)|>\Omega_h$, with $\dot\psi^G$ the GNSS course rate.
A candidate becomes an episode boundary when the quasi-static gravity directions on its two sides differ by more than $\theta_h$.
Orientation is estimated per episode $e$.

\paragraph{Gravity and tilt.}
Let $\bar{\mathbf f}^b$ and $\bar{\boldsymbol\omega}^b$ be low-passed at 0.5~Hz. The quasi-static set of episode $e$ is
\begin{equation}
\mathcal Q_e=\bigl\{n:\ \bigl|\lVert\bar{\mathbf f}^b[n]\rVert-g\bigr|<\epsilon_g,\ \lVert\bar{\boldsymbol\omega}^b[n]\rVert<\epsilon_\omega,\ |a^G_{\mathrm{lon}}(t_n)|<\epsilon_a,\ v(t_n)>v_{\min}\bigr\},
\end{equation}
and the up direction is $\mathbf u_e=\hat{\mathbf g}_e/\lVert\hat{\mathbf g}_e\rVert$ with $\hat{\mathbf g}_e$ the geometric median of $\{\bar{\mathbf f}^b[n]:n\in\mathcal Q_e\}$, which is robust to outliers and equivariant to rotation.
Gating matters: an ungated low-pass estimate absorbs sustained acceleration and tilts by $\arctan(a/g)$, i.e.\ $17^\circ$ at $3~\mathrm{m\,s^{-2}}$.
With $\mathbf k=\mathbf u_e\times\mathbf e_z/\lVert\mathbf u_e\times\mathbf e_z\rVert$ and $\theta=\operatorname{atan2}(\lVert\mathbf u_e\times\mathbf e_z\rVert,\mathbf u_e^\top\mathbf e_z)$, Rodrigues' formula gives the leveling rotation
\begin{equation}
\mathbf R_1=\mathbf I+\sin\theta\,[\mathbf k]_\times+(1-\cos\theta)\,[\mathbf k]_\times^2 .
\end{equation}
The route-average road grade biases this tilt (a 5\% grade is $2.9^\circ$).
The resulting error is at most 0.13\% on the vertical channel, but up to $0.5~\mathrm{m\,s^{-2}}$ on the longitudinal channel.
We therefore take longitudinal maneuver thresholds from GNSS rather than from the IMU.

\paragraph{Heading.}
After leveling, the horizontal axes differ from the vehicle axes by an unknown yaw $\psi$.
The leveled yaw rate $\omega^\ell_z=(\mathbf R_1\bar{\boldsymbol\omega}^b)_z$ does not depend on $\psi$.
This yields two GNSS-referenced targets at every fix: the longitudinal acceleration $a^G_{\mathrm{lon},k}$ and the centripetal acceleration $a^G_{\mathrm{lat},k}=v_k\,\omega^\ell_z(t_k)$.
With the leveled horizontal specific force $\mathbf h_k\in\R^2$ and $\mathbf r_k=(a^G_{\mathrm{lon},k},a^G_{\mathrm{lat},k})^\top$,
\begin{equation}
\hat\psi=\arg\min_{\psi}\sum_k\bigl\lVert\mathbf R_2(\psi)\mathbf h_k-\mathbf r_k\bigr\rVert^2
=\operatorname{atan2}\bigl(M_{21}-M_{12},\ M_{11}+M_{22}\bigr),\qquad \mathbf M=\sum_k\mathbf h_k\mathbf r_k^\top,
\label{eq:yaw}
\end{equation}
where $\mathbf R_2(\psi)=\bigl[\begin{smallmatrix}\cos\psi&\sin\psi\\-\sin\psi&\cos\psi\end{smallmatrix}\bigr]$.
The closed form follows from maximizing $\operatorname{tr}(\mathbf R_2\mathbf M)$.
Only fixes with $v_k>3~\mathrm{m\,s^{-1}}$ are used.
Finally $\mathbf R_{vb}=\mathbf R_z(\hat\psi)\mathbf R_1$, and a bootstrap over fixes provides the orientation confidence passed to Layer~3.

\paragraph{Speed decoupling.}
Vibration grows with speed.
For a road whose elevation PSD follows ISO~8608~\cite{iso8608}, $G_d(n)=G_d(n_0)(n/n_0)^{-w}$, the temporal input PSD at speed $v$ scales as $v^{w-1}f^{-w}$.
For a fixed vehicle response the RMS response therefore scales as $v^{(w-1)/2}$.
Rather than assume $w$, we estimate the elasticity with the window-level model $\log a_{w,i}=\alpha_{j(i)}+\gamma\log v_i+\varepsilon_i$, where $j$ indexes homogeneous segments, and report
\begin{equation}
\tilde a_w=a_w\,(v_{\mathrm{ref}}/\bar v)^{\hat\gamma},\qquad \bar v\ge v_{\min},\ v_{\mathrm{ref}}=10~\mathrm{m\,s^{-1}}.
\end{equation}
The raw $a_w$ describes how the ride felt and $\tilde a_w$ describes the road; narratives keep the two apart.

\paragraph{Frequency weighting.}
ISO~2631-1~\cite{iso2631} weights seated whole-body vibration with $W_k$ (vertical) and $W_d$ (horizontal).
Each is a product of band-limiting filters, an acceleration--velocity transition and, for $W_k$, an upward step ($p=j2\pi f$, $\omega_i=2\pi f_i$, $Q_1=1/\sqrt2$):
\begin{align}
H_h(p)&=\frac{p^2}{p^2+\omega_1p/Q_1+\omega_1^2}, &
H_l(p)&=\frac{\omega_2^2}{p^2+\omega_2p/Q_1+\omega_2^2}, \nonumber\\
H_t(p)&=\frac{(\omega_4^2/\omega_3)(p+\omega_3)}{p^2+\omega_4p/Q_4+\omega_4^2}, &
H_s(p)&=\frac{p^2+\omega_5p/Q_5+\omega_5^2}{p^2+\omega_6p/Q_6+\omega_6^2},
\end{align}
with $W=H_hH_lH_tH_s$ and the parameters of Table~\ref{tab:iso}.
Narratives are produced after the trip, so we apply $W$ in the frequency domain, $a_w=\mathcal F^{-1}\{\mathcal F\{a\}\,W(j2\pi f)\}$ with 10~s of zero padding.
This reproduces the analog magnitude and phase exactly below $f_s/2$.
For streaming use, a bilinear transform per second-order section, pre-warped at each section's corner frequency, is an alternative.
At $f_s=100$~Hz its magnitude error is at most 0.72~dB ($W_k$) and 1.22~dB ($W_d$) up to 20~Hz, but it reaches 7--8~dB at 40~Hz and underestimates impulsive peaks by 3.8\%.
At $f_s=400$~Hz the error stays within 0.23~dB up to 40~Hz.
The 100~Hz band limit $H_l$ cannot be realized for $f_s\le200$~Hz; omitting it changes $W$ by less than 0.11~dB below 40~Hz.
At $f_s=100$~Hz our analysis band ends near 40~Hz, below the 80~Hz upper limit of the standard.

\begin{table}[t]
\caption{ISO 2631-1 frequency-weighting parameters (Annex A). Frequencies in Hz.}
\label{tab:iso}
\begin{tabular}{lccccccccc}
\toprule
 & $f_1$ & $f_2$ & $f_3$ & $f_4$ & $Q_4$ & $f_5$ & $Q_5$ & $f_6$ & $Q_6$\\
\midrule
$W_k$ & 0.4 & 100 & 12.5 & 12.5 & 0.63 & 2.37 & 0.91 & 3.35 & 0.91\\
$W_d$ & 0.4 & 100 & 2.0 & 2.0 & 0.63 & -- & -- & -- & --\\
\bottomrule
\end{tabular}
\end{table}

\paragraph{Vibration measures.}
For each HRNU of duration $T$ and axis $j$ we compute
\begin{gather}
a_{w,j}=\Bigl(\tfrac1T\textstyle\int_0^Ta_{w,j}^2\,dt\Bigr)^{1/2},\qquad
a_v=\bigl(a_{w,x}^2+a_{w,y}^2+a_{w,z}^2\bigr)^{1/2},\\
\mathrm{MTVV}_j=\max_{t_0}\Bigl(\tfrac1\tau\textstyle\int_{t_0-\tau}^{t_0}a_{w,j}^2\,dt\Bigr)^{1/2}\ (\tau=1~\mathrm{s}),\qquad
\mathrm{VDV}_j=\Bigl(\textstyle\int_0^Ta_{w,j}^4\,dt\Bigr)^{1/4},\qquad
\mathrm{CF}_j=\frac{\max_t|a_{w,j}(t)|}{a_{w,j}}.
\end{gather}
Following the standard's evaluation rule, $a_w$ is always reported.
MTVV and VDV are added, and the unit is flagged as shock-dominated, when
\begin{equation}
S_j=\ind\bigl[\mathrm{CF}_j>9\ \lor\ \mathrm{MTVV}_j/a_{w,j}>1.5\ \lor\ \mathrm{VDV}_j/(a_{w,j}T^{1/4})>1.75\bigr].
\end{equation}
For a zero-mean stationary signal $\mathrm{VDV}/(a_wT^{1/4})=\kappa^{1/4}$, where $\kappa$ is the kurtosis.
The VDV criterion is therefore the condition $\kappa>9.38$; Gaussian vibration gives a ratio of 1.32.
Comfort descriptors follow the informative comfort scale of the standard (Annex~C).
Its bands overlap, so a single label uses geometric-mean cut-points $\boldsymbol\theta=(0.315,0.561,0.894,1.414,2.236)~\mathrm{m\,s^{-2}}$.
ISO~2631-1 defines comfort at the seat--occupant interface, whereas a phone in a holder measures the cabin structure.
We therefore call these quantities \emph{weighted cabin acceleration} and validate them against a phone fixed at the seat rail (Section~\ref{sec:data}).
Below about 1~Hz, $a_{w,x}$ and $a_{w,y}$ mostly reflect braking and cornering, and we report them as driving rather than road measures.

\paragraph{Modal band ratios.}
From the Welch PSD of $a^v_z$ (Hann, 5.12~s segments, 50\% overlap) we compute energies in a body band $[0.8,2.5]$~Hz, a wheel-hop band centred on the vehicle's wheel-hop peak $\hat f_{wh}\in[6,20]$~Hz, and a high band $[16,40]$~Hz.
From these we form $R_{wb}=\log(E_{\mathrm{wheel}}/E_{\mathrm{body}})$ and $R_{hb}=\log(E_{\mathrm{high}}/E_{\mathrm{body}})$.
Through $\lambda=v/f$ these bands correspond to wavelengths from about 0.25~m to 12.5~m at 10~m\,s$^{-1}$, i.e.\ unevenness and megatexture, not pavement macrotexture (0.5--50~mm~\cite{iso13473}).
They describe how the vehicle's modes are excited, and we use them within a vehicle and for segmentation.

\paragraph{Impulsive events and maneuver context.}
We band-pass the vertical channel, $\tilde z=\mathrm{BP}_{[4,25]\,\mathrm{Hz}}(a^v_z)$, and apply the Teager--Kaiser energy operator~\cite{kaiser1990tkeo},
\begin{equation}
\Psi[n]=\tilde z[n]^2-\tilde z[n-1]\,\tilde z[n+1],\qquad E[n]=\bigl(h_{50\,\mathrm{ms}}*\max(\Psi,0)\bigr)[n].
\end{equation}
For $\tilde z=A\cos(\Omega n+\phi)$, $\Psi=A^2\sin^2\Omega$, so the operator responds to both amplitude and frequency.
Candidates are maximal runs where $E$ exceeds $\theta[n]=\operatorname{med}_W(E)+\kappa\cdot1.4826\operatorname{MAD}_W(E)$ over a 10~s centred window.
The adaptive threshold follows the background, which is much higher on cobblestone than on asphalt.
Runs closer than 0.3~s are merged, and runs longer than $D_{\max}=0.5~\mathrm{s}+3~\mathrm{m}/\bar v$ are passed to segmentation as rough stretches.
The maneuver state $m_e$ is one of \textsc{brake}, \textsc{accel}, \textsc{turn}, \textsc{swerve}, \textsc{stop-go} or \textsc{cruise}.
Longitudinal states come from GNSS ($a^G_{\mathrm{lon}}$) and lateral ones from the yaw rate ($v\,\omega^v_z$); \textsc{stop-go} requires $v<2~\mathrm{m\,s^{-1}}$ and a longitudinal jerk $|j_x|>J_{sg}$.
With $\rho_e=E_{[8,25]}(e)/E_{[0.5,4]}(e)$ the band-energy ratio of $a^v_z$ in the event window, each candidate is labelled
\begin{equation}
\ell(e)=
\begin{cases}
\text{\textsc{artifact}} & m_e\in\{\text{\textsc{stop-go}},\text{\textsc{accel}}\}\ \land\ \max_{n\in e}|j_x[n]|>J_x\ \land\ \rho_e<\rho_0,\\
\text{\textsc{defect}} & \lnot\,\text{\textsc{artifact}}\ \land\ \rho_e\ge\rho_0\ \land\ d_e\le D_{\max},\\
\text{\textsc{ambiguous}} & \text{otherwise}.
\end{cases}
\label{eq:gate}
\end{equation}
Braking and turning never veto a defect.
In our data they do not increase vertical vibration (Section~\ref{sec:l1eval}), and drivers slow down \emph{before} speed bumps, so a braking veto would remove true events.
Acceleration is gated because it is the only maneuver that adds a dashboard-only high-band component.
A swerve without a vertical candidate within $\pm2$~s yields a separate low-confidence \emph{avoidance} event.
A logistic model trained on held-out routes maps the event features to $p_{\mathrm{defect}}$.

% ---------------------------------------------------------------------
\subsection{Layer 2: Spatial Grounding and Segmentation}
\label{sec:layer2}

\paragraph{Map matching.}
We match GNSS fixes to OpenStreetMap with a hidden Markov model~\cite{newson2009hmm}.
The states are candidate projections $x_{k,i}$ within 50~m of fix $z_k$, the emissions are $\mathcal N(\lVert z_k-x_{k,i}\rVert_{gc};0,\sigma_z^2)$, and the transitions are $\beta^{-1}\exp(-|\lVert z_k-z_{k+1}\rVert_{gc}-\lVert x_{k,i}-x_{k+1,j}\rVert_{\mathrm{route}}|/\beta)$, decoded with Viterbi.
Forward--backward marginals give a per-fix match confidence.
The matched arc length $s(t)$ replaces $\int v\,dt$.
Each matched way contributes its \texttt{surface}, \texttt{smoothness}, \texttt{highway} and \texttt{name} tags, and nodes such as \texttt{traffic\_calming} and level crossings become candidate anchors.

\paragraph{Map-informed change-point detection.}
We resample Layer-1 features on a 5~m along-track grid, $\mathbf y_j=[\log\tilde a_{w,z},\log a_{w,z},R_{wb},R_{hb},\text{event rate},\log a_{w,y}]_j$, robust-standardized per trip.
Speed is excluded so that boundaries do not follow speed changes.
Segmentation solves
\begin{equation}
\min_{m,\,\tau_{1:m}}\ \sum_{i=1}^{m+1}C\bigl(\mathbf y_{\tau_{i-1}+1:\tau_i}\bigr)+\sum_{i=1}^{m}\beta(\tau_i),\qquad
C(\mathbf y_{a:b})=\sum_{j=a}^{b}\lVert\mathbf y_j-\bar{\mathbf y}_{a:b}\rVert^2_{\mathbf D^{-1}},
\label{eq:pelt}
\end{equation}
with a penalty that falls near OpenStreetMap transitions $b\in\mathcal B$ (surface-tag changes, road-class changes, junctions):
\begin{equation}
\beta(\tau)=\beta_0\bigl(1-\lambda\,\pi(\tau)\bigr),\qquad
\pi(\tau)=\max_{b\in\mathcal B}w_b\exp\!\Bigl(-\frac{(s_\tau-s_b)^2}{2\sigma_b^2}\Bigr),\qquad \lambda\in[0,1).
\end{equation}
Setting $\lambda=0$ recovers PELT~\cite{killick2012pelt}.
Let $F(0)=0$, $G(t)=F(t)+\beta(t)\ind[t>0]$ and $F(s)=\min_{t\in\mathcal R_s}\{G(t)+C(\mathbf y_{t+1:s})\}$.

\begin{lemma}[Pruning under location-dependent penalties]
\label{lem:prune}
Suppose $C(\mathbf y_{t+1:s})+C(\mathbf y_{s+1:T})\le C(\mathbf y_{t+1:T})$ for all $t<s<T$.
If $G(t)+C(\mathbf y_{t+1:s})\ge G(s)$, then $t$ is not the optimal last change point for any $T>s$ and can be removed from $\mathcal R$.
\end{lemma}
\begin{proof}
For $T>s$, $G(t)+C(\mathbf y_{t+1:T})\ge G(t)+C(\mathbf y_{t+1:s})+C(\mathbf y_{s+1:T})\ge G(s)+C(\mathbf y_{s+1:T})$.
\end{proof}

The usual condition $F(t)+C(\mathbf y_{t+1:s})\ge F(s)$ coincides with Lemma~\ref{lem:prune} only for a constant penalty.
With a lower penalty at map transitions it discards valid candidates.
In 300 random series it returned a sub-optimal partition every time, whereas the rule of Lemma~\ref{lem:prune} matched exhaustive optimal partitioning in all 300.
We choose $\beta_0=c\,d\ln n$ and $(c,\lambda)$ by leave-one-route-out validation, and report the penalty path~\cite{haynes2017crops}.

\paragraph{Arbitration and HRNUs.}
A change point within 25~m of a relevant map transition is snapped to it: the map fixes the location and the signal decides existence.
A map transition without a change point is kept only if a two-sample Hotelling test on $\mathbf y$ across it rejects at $\alpha=0.01$; otherwise every way split in the map would start a new unit.
Units shorter than 50~m are merged into their closer neighbour.
Each boundary is typed as map+signal, signal-only or map-only.
Each unit carries a map--signal disagreement score $1-P(\text{surface}=\text{OSM tag}\mid\mathbf y)$ when the map has a surface tag.
The resulting HRNUs, with their measures, events, anchors and confidences, are the input to Layer~3.

% ---------------------------------------------------------------------
\subsection{Layer 3: Kinetic--Spatial Tokenization and Grounded Generation}
\label{sec:layer3}

\paragraph{Claim ledger.}
Layer~3 first converts HRNUs and events into a ledger $\mathcal L$ of typed facts with identifiers, values, tolerances and confidences.
The types are severity, numeric, event, anchor, comparison, missing and metadata.
A \textsc{missing} fact records that a value is unavailable and why, e.g.\ a 48~s GNSS gap.
The ledger is the reference for both the run-time verifier and the evaluation.

\paragraph{Kinetic--Spatial Tokenizer.}
Algorithm~\ref{alg:kst} compiles the ledger into a capsule that fits a token budget $B$.
Each unit has a salience $\sigma_m$ and three serializations (omit, brief, full) with token costs $c_{m,\ell}$ measured with the target tokenizer.
Selection solves a multiple-choice knapsack exactly by dynamic programming:
\begin{equation}
V_m(b)=\max_{\ell\in\Lambda_m}\bigl\{V_{m-1}(b-c_{m,\ell})+\sigma_m\eta_\ell\bigr\},\qquad V_0(b)=0,
\label{eq:knap}
\end{equation}
where $\Lambda_m$ excludes \textsc{omit} for units in the must-mention set $\mathcal M$ and $\eta_{\textsc{omit}}=0$.
The cost is $O(3MB)$ for $M$ units.

\begin{algorithm}[t]
\caption{Kinetic--Spatial Tokenizer (KST)}
\label{alg:kst}
\begin{algorithmic}[1]
\Require HRNUs $\mathcal U=(u_1,\dots,u_M)$ in route order; events $\mathcal E$; anchors $\mathcal A$; optional history $\mathcal H$; audience $\pi$; budget $B$; tokenizer $\mathrm{tok}$
\Ensure capsule $\kappa$ with $|\mathrm{tok}(\kappa)|\le B$; ledger $\mathcal L$
\State $\mathcal L\gets\textsc{MetaFacts}(\text{channels},\text{map coverage},\pi)$
\For{$m=1,\dots,M$}
  \State $k_m\gets\textsc{IsoClass}(a_{v,m})$;\quad $\tilde k_m\gets\textsc{IsoClass}(\tilde a_{w,m})$ \Comment{felt vs.\ road}
  \State $\mathcal E_m\gets\{e\in\mathcal E: s(e)\in[s^0_m,s^1_m)\}$
  \State $\alpha_m\gets$ nearest anchor in $\mathcal A$ within $r_{\max}$ of $u_m$, or $\varnothing$
  \State $\sigma_m\gets w_1k_m+w_2\sum_{e\in\mathcal E_m}p_e\,\mathrm{sev}(e)+w_3\,\Delta_{\mathcal H}(u_m)+w_4\,\delta^{\mathrm{map}}_m$
  \State $\mathcal L\gets\mathcal L\cup\textsc{Facts}(u_m,k_m,\tilde k_m,\mathcal E_m,\alpha_m)$ \Comment{including \textsc{missing} facts}
  \For{$\ell\in\{\textsc{omit},\textsc{brief},\textsc{full}\}$}
    \State $c_{m,\ell}\gets|\mathrm{tok}(\textsc{Ser}_\ell(u_m))|$
  \EndFor
\EndFor
\State $\mathcal M\gets\{m:\exists e\in\mathcal E_m,\ p_e\,\mathrm{sev}(e)\ge\eta^\star\}$ \Comment{must-mention units}
\State $B'\gets B-|\mathrm{tok}(\textsc{Header})|-c_{\mathrm{rest}}$
\State $\mathbf x^\star\gets\textsc{McKnapsack}(\{c_{m,\ell}\},\{\sigma_m\eta_\ell\},B',\mathcal M)$ \Comment{Eq.~\eqref{eq:knap}}
\If{$\mathbf x^\star$ is infeasible}
  \State merge adjacent units outside $\mathcal M$ and repeat; if still infeasible, increase $B$ and log it
\EndIf
\State $\kappa\gets\textsc{Header}\ \|\ \textsc{Ser}_{x^\star_1}(u_1)\ \|\cdots\|\ \textsc{Ser}_{x^\star_M}(u_M)\ \|\ \textsc{Rest}(\{u_m:x^\star_m=\textsc{omit}\})$
\State \Return $\kappa,\ \mathcal L$
\end{algorithmic}
\end{algorithm}

\paragraph{Plan--ground--realize.}
Generation runs in four steps.
\emph{Plan:} the SLM emits a JSON plan listing which units and facts to mention, their order, a rhetorical relation for each (elaboration, contrast, cause, comparison with history) and a hedge.
Decoding is constrained by a grammar generated per capsule from the ledger identifiers~\cite{ggerganov_llamacpp}, so the plan is always well-formed and cannot reference a unit or event that does not exist.
\emph{Ground:} deterministic code expands each plan item into a fact frame with values, names and units taken from $\mathcal L$.
\emph{Realize:} the SLM writes text from the fact frames, either with placeholders for numbers and names that are filled afterwards, or as free text.
\emph{Verify:} claims extracted from the text (Section~\ref{sec:metrics}) are checked against $\mathcal L$.
A failing sentence is regenerated at most twice and then replaced by the template realization of the same plan item.
The SLM therefore controls selection, ordering, aggregation and phrasing, but not values.

\paragraph{Hedging.}
The plan assigns hedges from fact confidence $p$: assertive for $p\ge0.9$, \emph{likely} for $0.7\le p<0.9$, \emph{possibly} for $0.4\le p<0.7$; facts with $p<0.4$ are mentioned only when asked about.
The verifier rejects a sentence whose hedge is stronger than its confidence band allows.
 inside an acmart document.
% Preamble requirements (acmart already loads amsmath, amsthm, booktabs, xcolor):
%   \usepackage{algorithm}
%   \usepackage{algpseudocode}
% Bibliography keys used below (add to your .bib):
%   iso2631, iso8855, iso8608, iso13473, newson2009hmm, killick2012pelt,
%   haynes2017crops, kaiser1990tkeo, mohan2008nericell, eriksson2008pothole,
%   menegazzo2021pvs, ggerganov_llamacpp
% =====================================================================
\providecommand{\R}{\mathbb{R}}
\providecommand{\ind}{\mathbf{1}}
\providecommand{\TBD}[1][]{\textcolor{red}{[TBD#1]}}

\section{The RoadScribe System}
\label{sec:system}

RoadScribe turns a trip recorded by a commodity smartphone into a short narrative and answers follow-up questions about it.
The pipeline has three layers (Figure~\ref{fig:arch}).
Layer~1 converts inertial and GNSS samples into vehicle-frame ride measures and candidate events.
Layer~2 anchors these measures to the road network and partitions the trip into homogeneous road narrative units (HRNUs).
Layer~3 compiles the units into a token-budgeted capsule and a machine-checkable claim ledger, and generates text with an on-device small language model (SLM) under a plan--ground--realize protocol.

Four principles shape the design.
(P1)~Every number, place name and event in the output can be traced to a ledger fact.
(P2)~The language model decides what to say and how to say it, but never supplies a value; values come from the ledger.
(P3)~Uncertainty estimated in Layers~1--2 reaches the text as calibrated hedges.
(P4)~When verification fails, a deterministic realizer takes over, so the SLM can add coverage but cannot remove faithfulness.

\begin{figure}[t]
  \centering
  \fbox{\parbox{0.95\linewidth}{\centering\vspace{2em}Architecture figure placeholder: IMU+GNSS $\rightarrow$ Layer 1 (orientation, ISO weighting, events) $\rightarrow$ Layer 2 (map matching, MI-PELT, HRNUs) $\rightarrow$ Layer 3 (ledger, KST capsule, plan--ground--realize, verifier, template fallback) $\rightarrow$ narrative / answers\vspace{2em}}}
  \caption{RoadScribe architecture.}
  \label{fig:arch}
\end{figure}

% ---------------------------------------------------------------------
\subsection{Layer 1: Kinetic Signal Processing}
\label{sec:layer1}

\paragraph{Notation.}
The IMU delivers specific force $\mathbf f^b[n]\in\R^3$ and angular rate $\boldsymbol\omega^b[n]\in\R^3$ in the phone body frame $b$ at rate $f_s$ ($T_s=1/f_s$).
GNSS delivers speed $v_k$ at times $t_k$ (1~Hz).
The vehicle frame $v$ follows ISO~8855~\cite{iso8855} ($x$ forward, $y$ left, $z$ up).
We estimate $\mathbf R_{vb}\in SO(3)$ and the gravity-free acceleration
\begin{equation}
  \mathbf a^v[n]=\mathbf R_{vb}\,\mathbf f^b[n]-g\,\mathbf e_z .
  \label{eq:av}
\end{equation}
Speed is interpolated to IMU time with a monotone cubic (PCHIP), giving $v(t)$ and $a^G_{\mathrm{lon}}(t)=\dot v(t)$.

\paragraph{Mount episodes.}
A phone in a holder is static relative to the vehicle except when handled.
We mark candidate handling where the gyroscope reports rotation that the vehicle does not perform, $\lVert\bar{\boldsymbol\omega}^b[n]\rVert-|\dot\psi^G(t_n)|>\Omega_h$, with $\dot\psi^G$ the GNSS course rate.
A candidate becomes an episode boundary when the quasi-static gravity directions on its two sides differ by more than $\theta_h$.
Orientation is estimated per episode $e$.

\paragraph{Gravity and tilt.}
Let $\bar{\mathbf f}^b$ and $\bar{\boldsymbol\omega}^b$ be low-passed at 0.5~Hz. The quasi-static set of episode $e$ is
\begin{equation}
\mathcal Q_e=\bigl\{n:\ \bigl|\lVert\bar{\mathbf f}^b[n]\rVert-g\bigr|<\epsilon_g,\ \lVert\bar{\boldsymbol\omega}^b[n]\rVert<\epsilon_\omega,\ |a^G_{\mathrm{lon}}(t_n)|<\epsilon_a,\ v(t_n)>v_{\min}\bigr\},
\end{equation}
and the up direction is $\mathbf u_e=\hat{\mathbf g}_e/\lVert\hat{\mathbf g}_e\rVert$ with $\hat{\mathbf g}_e$ the geometric median of $\{\bar{\mathbf f}^b[n]:n\in\mathcal Q_e\}$, which is robust to outliers and equivariant to rotation.
Gating matters: an ungated low-pass estimate absorbs sustained acceleration and tilts by $\arctan(a/g)$, i.e.\ $17^\circ$ at $3~\mathrm{m\,s^{-2}}$.
With $\mathbf k=\mathbf u_e\times\mathbf e_z/\lVert\mathbf u_e\times\mathbf e_z\rVert$ and $\theta=\operatorname{atan2}(\lVert\mathbf u_e\times\mathbf e_z\rVert,\mathbf u_e^\top\mathbf e_z)$, Rodrigues' formula gives the leveling rotation
\begin{equation}
\mathbf R_1=\mathbf I+\sin\theta\,[\mathbf k]_\times+(1-\cos\theta)\,[\mathbf k]_\times^2 .
\end{equation}
The route-average road grade biases this tilt (a 5\% grade is $2.9^\circ$).
The resulting error is at most 0.13\% on the vertical channel, but up to $0.5~\mathrm{m\,s^{-2}}$ on the longitudinal channel.
We therefore take longitudinal maneuver thresholds from GNSS rather than from the IMU.

\paragraph{Heading.}
After leveling, the horizontal axes differ from the vehicle axes by an unknown yaw $\psi$.
The leveled yaw rate $\omega^\ell_z=(\mathbf R_1\bar{\boldsymbol\omega}^b)_z$ does not depend on $\psi$.
This yields two GNSS-referenced targets at every fix: the longitudinal acceleration $a^G_{\mathrm{lon},k}$ and the centripetal acceleration $a^G_{\mathrm{lat},k}=v_k\,\omega^\ell_z(t_k)$.
With the leveled horizontal specific force $\mathbf h_k\in\R^2$ and $\mathbf r_k=(a^G_{\mathrm{lon},k},a^G_{\mathrm{lat},k})^\top$,
\begin{equation}
\hat\psi=\arg\min_{\psi}\sum_k\bigl\lVert\mathbf R_2(\psi)\mathbf h_k-\mathbf r_k\bigr\rVert^2
=\operatorname{atan2}\bigl(M_{21}-M_{12},\ M_{11}+M_{22}\bigr),\qquad \mathbf M=\sum_k\mathbf h_k\mathbf r_k^\top,
\label{eq:yaw}
\end{equation}
where $\mathbf R_2(\psi)=\bigl[\begin{smallmatrix}\cos\psi&\sin\psi\\-\sin\psi&\cos\psi\end{smallmatrix}\bigr]$.
The closed form follows from maximizing $\operatorname{tr}(\mathbf R_2\mathbf M)$.
Only fixes with $v_k>3~\mathrm{m\,s^{-1}}$ are used.
Finally $\mathbf R_{vb}=\mathbf R_z(\hat\psi)\mathbf R_1$, and a bootstrap over fixes provides the orientation confidence passed to Layer~3.

\paragraph{Speed decoupling.}
Vibration grows with speed.
For a road whose elevation PSD follows ISO~8608~\cite{iso8608}, $G_d(n)=G_d(n_0)(n/n_0)^{-w}$, the temporal input PSD at speed $v$ scales as $v^{w-1}f^{-w}$.
For a fixed vehicle response the RMS response therefore scales as $v^{(w-1)/2}$.
Rather than assume $w$, we estimate the elasticity with the window-level model $\log a_{w,i}=\alpha_{j(i)}+\gamma\log v_i+\varepsilon_i$, where $j$ indexes homogeneous segments, and report
\begin{equation}
\tilde a_w=a_w\,(v_{\mathrm{ref}}/\bar v)^{\hat\gamma},\qquad \bar v\ge v_{\min},\ v_{\mathrm{ref}}=10~\mathrm{m\,s^{-1}}.
\end{equation}
The raw $a_w$ describes how the ride felt and $\tilde a_w$ describes the road; narratives keep the two apart.

\paragraph{Frequency weighting.}
ISO~2631-1~\cite{iso2631} weights seated whole-body vibration with $W_k$ (vertical) and $W_d$ (horizontal).
Each is a product of band-limiting filters, an acceleration--velocity transition and, for $W_k$, an upward step ($p=j2\pi f$, $\omega_i=2\pi f_i$, $Q_1=1/\sqrt2$):
\begin{align}
H_h(p)&=\frac{p^2}{p^2+\omega_1p/Q_1+\omega_1^2}, &
H_l(p)&=\frac{\omega_2^2}{p^2+\omega_2p/Q_1+\omega_2^2}, \nonumber\\
H_t(p)&=\frac{(\omega_4^2/\omega_3)(p+\omega_3)}{p^2+\omega_4p/Q_4+\omega_4^2}, &
H_s(p)&=\frac{p^2+\omega_5p/Q_5+\omega_5^2}{p^2+\omega_6p/Q_6+\omega_6^2},
\end{align}
with $W=H_hH_lH_tH_s$ and the parameters of Table~\ref{tab:iso}.
Narratives are produced after the trip, so we apply $W$ in the frequency domain, $a_w=\mathcal F^{-1}\{\mathcal F\{a\}\,W(j2\pi f)\}$ with 10~s of zero padding.
This reproduces the analog magnitude and phase exactly below $f_s/2$.
For streaming use, a bilinear transform per second-order section, pre-warped at each section's corner frequency, is an alternative.
At $f_s=100$~Hz its magnitude error is at most 0.72~dB ($W_k$) and 1.22~dB ($W_d$) up to 20~Hz, but it reaches 7--8~dB at 40~Hz and underestimates impulsive peaks by 3.8\%.
At $f_s=400$~Hz the error stays within 0.23~dB up to 40~Hz.
The 100~Hz band limit $H_l$ cannot be realized for $f_s\le200$~Hz; omitting it changes $W$ by less than 0.11~dB below 40~Hz.
At $f_s=100$~Hz our analysis band ends near 40~Hz, below the 80~Hz upper limit of the standard.

\begin{table}[t]
\caption{ISO 2631-1 frequency-weighting parameters (Annex A). Frequencies in Hz.}
\label{tab:iso}
\begin{tabular}{lccccccccc}
\toprule
 & $f_1$ & $f_2$ & $f_3$ & $f_4$ & $Q_4$ & $f_5$ & $Q_5$ & $f_6$ & $Q_6$\\
\midrule
$W_k$ & 0.4 & 100 & 12.5 & 12.5 & 0.63 & 2.37 & 0.91 & 3.35 & 0.91\\
$W_d$ & 0.4 & 100 & 2.0 & 2.0 & 0.63 & -- & -- & -- & --\\
\bottomrule
\end{tabular}
\end{table}

\paragraph{Vibration measures.}
For each HRNU of duration $T$ and axis $j$ we compute
\begin{gather}
a_{w,j}=\Bigl(\tfrac1T\textstyle\int_0^Ta_{w,j}^2\,dt\Bigr)^{1/2},\qquad
a_v=\bigl(a_{w,x}^2+a_{w,y}^2+a_{w,z}^2\bigr)^{1/2},\\
\mathrm{MTVV}_j=\max_{t_0}\Bigl(\tfrac1\tau\textstyle\int_{t_0-\tau}^{t_0}a_{w,j}^2\,dt\Bigr)^{1/2}\ (\tau=1~\mathrm{s}),\qquad
\mathrm{VDV}_j=\Bigl(\textstyle\int_0^Ta_{w,j}^4\,dt\Bigr)^{1/4},\qquad
\mathrm{CF}_j=\frac{\max_t|a_{w,j}(t)|}{a_{w,j}}.
\end{gather}
Following the standard's evaluation rule, $a_w$ is always reported.
MTVV and VDV are added, and the unit is flagged as shock-dominated, when
\begin{equation}
S_j=\ind\bigl[\mathrm{CF}_j>9\ \lor\ \mathrm{MTVV}_j/a_{w,j}>1.5\ \lor\ \mathrm{VDV}_j/(a_{w,j}T^{1/4})>1.75\bigr].
\end{equation}
For a zero-mean stationary signal $\mathrm{VDV}/(a_wT^{1/4})=\kappa^{1/4}$, where $\kappa$ is the kurtosis.
The VDV criterion is therefore the condition $\kappa>9.38$; Gaussian vibration gives a ratio of 1.32.
Comfort descriptors follow the informative comfort scale of the standard (Annex~C).
Its bands overlap, so a single label uses geometric-mean cut-points $\boldsymbol\theta=(0.315,0.561,0.894,1.414,2.236)~\mathrm{m\,s^{-2}}$.
ISO~2631-1 defines comfort at the seat--occupant interface, whereas a phone in a holder measures the cabin structure.
We therefore call these quantities \emph{weighted cabin acceleration} and validate them against a phone fixed at the seat rail (Section~\ref{sec:data}).
Below about 1~Hz, $a_{w,x}$ and $a_{w,y}$ mostly reflect braking and cornering, and we report them as driving rather than road measures.

\paragraph{Modal band ratios.}
From the Welch PSD of $a^v_z$ (Hann, 5.12~s segments, 50\% overlap) we compute energies in a body band $[0.8,2.5]$~Hz, a wheel-hop band centred on the vehicle's wheel-hop peak $\hat f_{wh}\in[6,20]$~Hz, and a high band $[16,40]$~Hz.
From these we form $R_{wb}=\log(E_{\mathrm{wheel}}/E_{\mathrm{body}})$ and $R_{hb}=\log(E_{\mathrm{high}}/E_{\mathrm{body}})$.
Through $\lambda=v/f$ these bands correspond to wavelengths from about 0.25~m to 12.5~m at 10~m\,s$^{-1}$, i.e.\ unevenness and megatexture, not pavement macrotexture (0.5--50~mm~\cite{iso13473}).
They describe how the vehicle's modes are excited, and we use them within a vehicle and for segmentation.

\paragraph{Impulsive events and maneuver context.}
We band-pass the vertical channel, $\tilde z=\mathrm{BP}_{[4,25]\,\mathrm{Hz}}(a^v_z)$, and apply the Teager--Kaiser energy operator~\cite{kaiser1990tkeo},
\begin{equation}
\Psi[n]=\tilde z[n]^2-\tilde z[n-1]\,\tilde z[n+1],\qquad E[n]=\bigl(h_{50\,\mathrm{ms}}*\max(\Psi,0)\bigr)[n].
\end{equation}
For $\tilde z=A\cos(\Omega n+\phi)$, $\Psi=A^2\sin^2\Omega$, so the operator responds to both amplitude and frequency.
Candidates are maximal runs where $E$ exceeds $\theta[n]=\operatorname{med}_W(E)+\kappa\cdot1.4826\operatorname{MAD}_W(E)$ over a 10~s centred window.
The adaptive threshold follows the background, which is much higher on cobblestone than on asphalt.
Runs closer than 0.3~s are merged, and runs longer than $D_{\max}=0.5~\mathrm{s}+3~\mathrm{m}/\bar v$ are passed to segmentation as rough stretches.
The maneuver state $m_e$ is one of \textsc{brake}, \textsc{accel}, \textsc{turn}, \textsc{swerve}, \textsc{stop-go} or \textsc{cruise}.
Longitudinal states come from GNSS ($a^G_{\mathrm{lon}}$) and lateral ones from the yaw rate ($v\,\omega^v_z$); \textsc{stop-go} requires $v<2~\mathrm{m\,s^{-1}}$ and a longitudinal jerk $|j_x|>J_{sg}$.
With $\rho_e=E_{[8,25]}(e)/E_{[0.5,4]}(e)$ the band-energy ratio of $a^v_z$ in the event window, each candidate is labelled
\begin{equation}
\ell(e)=
\begin{cases}
\text{\textsc{artifact}} & m_e\in\{\text{\textsc{stop-go}},\text{\textsc{accel}}\}\ \land\ \max_{n\in e}|j_x[n]|>J_x\ \land\ \rho_e<\rho_0,\\
\text{\textsc{defect}} & \lnot\,\text{\textsc{artifact}}\ \land\ \rho_e\ge\rho_0\ \land\ d_e\le D_{\max},\\
\text{\textsc{ambiguous}} & \text{otherwise}.
\end{cases}
\label{eq:gate}
\end{equation}
Braking and turning never veto a defect.
In our data they do not increase vertical vibration (Section~\ref{sec:l1eval}), and drivers slow down \emph{before} speed bumps, so a braking veto would remove true events.
Acceleration is gated because it is the only maneuver that adds a dashboard-only high-band component.
A swerve without a vertical candidate within $\pm2$~s yields a separate low-confidence \emph{avoidance} event.
A logistic model trained on held-out routes maps the event features to $p_{\mathrm{defect}}$.

% ---------------------------------------------------------------------
\subsection{Layer 2: Spatial Grounding and Segmentation}
\label{sec:layer2}

\paragraph{Map matching.}
We match GNSS fixes to OpenStreetMap with a hidden Markov model~\cite{newson2009hmm}.
The states are candidate projections $x_{k,i}$ within 50~m of fix $z_k$, the emissions are $\mathcal N(\lVert z_k-x_{k,i}\rVert_{gc};0,\sigma_z^2)$, and the transitions are $\beta^{-1}\exp(-|\lVert z_k-z_{k+1}\rVert_{gc}-\lVert x_{k,i}-x_{k+1,j}\rVert_{\mathrm{route}}|/\beta)$, decoded with Viterbi.
Forward--backward marginals give a per-fix match confidence.
The matched arc length $s(t)$ replaces $\int v\,dt$.
Each matched way contributes its \texttt{surface}, \texttt{smoothness}, \texttt{highway} and \texttt{name} tags, and nodes such as \texttt{traffic\_calming} and level crossings become candidate anchors.

\paragraph{Map-informed change-point detection.}
We resample Layer-1 features on a 5~m along-track grid, $\mathbf y_j=[\log\tilde a_{w,z},\log a_{w,z},R_{wb},R_{hb},\text{event rate},\log a_{w,y}]_j$, robust-standardized per trip.
Speed is excluded so that boundaries do not follow speed changes.
Segmentation solves
\begin{equation}
\min_{m,\,\tau_{1:m}}\ \sum_{i=1}^{m+1}C\bigl(\mathbf y_{\tau_{i-1}+1:\tau_i}\bigr)+\sum_{i=1}^{m}\beta(\tau_i),\qquad
C(\mathbf y_{a:b})=\sum_{j=a}^{b}\lVert\mathbf y_j-\bar{\mathbf y}_{a:b}\rVert^2_{\mathbf D^{-1}},
\label{eq:pelt}
\end{equation}
with a penalty that falls near OpenStreetMap transitions $b\in\mathcal B$ (surface-tag changes, road-class changes, junctions):
\begin{equation}
\beta(\tau)=\beta_0\bigl(1-\lambda\,\pi(\tau)\bigr),\qquad
\pi(\tau)=\max_{b\in\mathcal B}w_b\exp\!\Bigl(-\frac{(s_\tau-s_b)^2}{2\sigma_b^2}\Bigr),\qquad \lambda\in[0,1).
\end{equation}
Setting $\lambda=0$ recovers PELT~\cite{killick2012pelt}.
Let $F(0)=0$, $G(t)=F(t)+\beta(t)\ind[t>0]$ and $F(s)=\min_{t\in\mathcal R_s}\{G(t)+C(\mathbf y_{t+1:s})\}$.

\begin{lemma}[Pruning under location-dependent penalties]
\label{lem:prune}
Suppose $C(\mathbf y_{t+1:s})+C(\mathbf y_{s+1:T})\le C(\mathbf y_{t+1:T})$ for all $t<s<T$.
If $G(t)+C(\mathbf y_{t+1:s})\ge G(s)$, then $t$ is not the optimal last change point for any $T>s$ and can be removed from $\mathcal R$.
\end{lemma}
\begin{proof}
For $T>s$, $G(t)+C(\mathbf y_{t+1:T})\ge G(t)+C(\mathbf y_{t+1:s})+C(\mathbf y_{s+1:T})\ge G(s)+C(\mathbf y_{s+1:T})$.
\end{proof}

The usual condition $F(t)+C(\mathbf y_{t+1:s})\ge F(s)$ coincides with Lemma~\ref{lem:prune} only for a constant penalty.
With a lower penalty at map transitions it discards valid candidates.
In 300 random series it returned a sub-optimal partition every time, whereas the rule of Lemma~\ref{lem:prune} matched exhaustive optimal partitioning in all 300.
We choose $\beta_0=c\,d\ln n$ and $(c,\lambda)$ by leave-one-route-out validation, and report the penalty path~\cite{haynes2017crops}.

\paragraph{Arbitration and HRNUs.}
A change point within 25~m of a relevant map transition is snapped to it: the map fixes the location and the signal decides existence.
A map transition without a change point is kept only if a two-sample Hotelling test on $\mathbf y$ across it rejects at $\alpha=0.01$; otherwise every way split in the map would start a new unit.
Units shorter than 50~m are merged into their closer neighbour.
Each boundary is typed as map+signal, signal-only or map-only.
Each unit carries a map--signal disagreement score $1-P(\text{surface}=\text{OSM tag}\mid\mathbf y)$ when the map has a surface tag.
The resulting HRNUs, with their measures, events, anchors and confidences, are the input to Layer~3.

% ---------------------------------------------------------------------
\subsection{Layer 3: Kinetic--Spatial Tokenization and Grounded Generation}
\label{sec:layer3}

\paragraph{Claim ledger.}
Layer~3 first converts HRNUs and events into a ledger $\mathcal L$ of typed facts with identifiers, values, tolerances and confidences.
The types are severity, numeric, event, anchor, comparison, missing and metadata.
A \textsc{missing} fact records that a value is unavailable and why, e.g.\ a 48~s GNSS gap.
The ledger is the reference for both the run-time verifier and the evaluation.

\paragraph{Kinetic--Spatial Tokenizer.}
Algorithm~\ref{alg:kst} compiles the ledger into a capsule that fits a token budget $B$.
Each unit has a salience $\sigma_m$ and three serializations (omit, brief, full) with token costs $c_{m,\ell}$ measured with the target tokenizer.
Selection solves a multiple-choice knapsack exactly by dynamic programming:
\begin{equation}
V_m(b)=\max_{\ell\in\Lambda_m}\bigl\{V_{m-1}(b-c_{m,\ell})+\sigma_m\eta_\ell\bigr\},\qquad V_0(b)=0,
\label{eq:knap}
\end{equation}
where $\Lambda_m$ excludes \textsc{omit} for units in the must-mention set $\mathcal M$ and $\eta_{\textsc{omit}}=0$.
The cost is $O(3MB)$ for $M$ units.

\begin{algorithm}[t]
\caption{Kinetic--Spatial Tokenizer (KST)}
\label{alg:kst}
\begin{algorithmic}[1]
\Require HRNUs $\mathcal U=(u_1,\dots,u_M)$ in route order; events $\mathcal E$; anchors $\mathcal A$; optional history $\mathcal H$; audience $\pi$; budget $B$; tokenizer $\mathrm{tok}$
\Ensure capsule $\kappa$ with $|\mathrm{tok}(\kappa)|\le B$; ledger $\mathcal L$
\State $\mathcal L\gets\textsc{MetaFacts}(\text{channels},\text{map coverage},\pi)$
\For{$m=1,\dots,M$}
  \State $k_m\gets\textsc{IsoClass}(a_{v,m})$;\quad $\tilde k_m\gets\textsc{IsoClass}(\tilde a_{w,m})$ \Comment{felt vs.\ road}
  \State $\mathcal E_m\gets\{e\in\mathcal E: s(e)\in[s^0_m,s^1_m)\}$
  \State $\alpha_m\gets$ nearest anchor in $\mathcal A$ within $r_{\max}$ of $u_m$, or $\varnothing$
  \State $\sigma_m\gets w_1k_m+w_2\sum_{e\in\mathcal E_m}p_e\,\mathrm{sev}(e)+w_3\,\Delta_{\mathcal H}(u_m)+w_4\,\delta^{\mathrm{map}}_m$
  \State $\mathcal L\gets\mathcal L\cup\textsc{Facts}(u_m,k_m,\tilde k_m,\mathcal E_m,\alpha_m)$ \Comment{including \textsc{missing} facts}
  \For{$\ell\in\{\textsc{omit},\textsc{brief},\textsc{full}\}$}
    \State $c_{m,\ell}\gets|\mathrm{tok}(\textsc{Ser}_\ell(u_m))|$
  \EndFor
\EndFor
\State $\mathcal M\gets\{m:\exists e\in\mathcal E_m,\ p_e\,\mathrm{sev}(e)\ge\eta^\star\}$ \Comment{must-mention units}
\State $B'\gets B-|\mathrm{tok}(\textsc{Header})|-c_{\mathrm{rest}}$
\State $\mathbf x^\star\gets\textsc{McKnapsack}(\{c_{m,\ell}\},\{\sigma_m\eta_\ell\},B',\mathcal M)$ \Comment{Eq.~\eqref{eq:knap}}
\If{$\mathbf x^\star$ is infeasible}
  \State merge adjacent units outside $\mathcal M$ and repeat; if still infeasible, increase $B$ and log it
\EndIf
\State $\kappa\gets\textsc{Header}\ \|\ \textsc{Ser}_{x^\star_1}(u_1)\ \|\cdots\|\ \textsc{Ser}_{x^\star_M}(u_M)\ \|\ \textsc{Rest}(\{u_m:x^\star_m=\textsc{omit}\})$
\State \Return $\kappa,\ \mathcal L$
\end{algorithmic}
\end{algorithm}

\paragraph{Plan--ground--realize.}
Generation runs in four steps.
\emph{Plan:} the SLM emits a JSON plan listing which units and facts to mention, their order, a rhetorical relation for each (elaboration, contrast, cause, comparison with history) and a hedge.
Decoding is constrained by a grammar generated per capsule from the ledger identifiers~\cite{ggerganov_llamacpp}, so the plan is always well-formed and cannot reference a unit or event that does not exist.
\emph{Ground:} deterministic code expands each plan item into a fact frame with values, names and units taken from $\mathcal L$.
\emph{Realize:} the SLM writes text from the fact frames, either with placeholders for numbers and names that are filled afterwards, or as free text.
\emph{Verify:} claims extracted from the text (Section~\ref{sec:metrics}) are checked against $\mathcal L$.
A failing sentence is regenerated at most twice and then replaced by the template realization of the same plan item.
The SLM therefore controls selection, ordering, aggregation and phrasing, but not values.

\paragraph{Hedging.}
The plan assigns hedges from fact confidence $p$: assertive for $p\ge0.9$, \emph{likely} for $0.7\le p<0.9$, \emph{possibly} for $0.4\le p<0.7$; facts with $p<0.4$ are mentioned only when asked about.
The verifier rejects a sentence whose hedge is stronger than its confidence band allows.
 inside an acmart document.
% Preamble requirements (acmart already loads amsmath, amsthm, booktabs, xcolor):
%   \usepackage{algorithm}
%   \usepackage{algpseudocode}
% Bibliography keys used below (add to your .bib):
%   iso2631, iso8855, iso8608, iso13473, newson2009hmm, killick2012pelt,
%   haynes2017crops, kaiser1990tkeo, mohan2008nericell, eriksson2008pothole,
%   menegazzo2021pvs, ggerganov_llamacpp
% =====================================================================
\providecommand{\R}{\mathbb{R}}
\providecommand{\ind}{\mathbf{1}}
\providecommand{\TBD}[1][]{\textcolor{red}{[TBD#1]}}

\section{The RoadScribe System}
\label{sec:system}

RoadScribe turns a trip recorded by a commodity smartphone into a short narrative and answers follow-up questions about it.
The pipeline has three layers (Figure~\ref{fig:arch}).
Layer~1 converts inertial and GNSS samples into vehicle-frame ride measures and candidate events.
Layer~2 anchors these measures to the road network and partitions the trip into homogeneous road narrative units (HRNUs).
Layer~3 compiles the units into a token-budgeted capsule and a machine-checkable claim ledger, and generates text with an on-device small language model (SLM) under a plan--ground--realize protocol.

Four principles shape the design.
(P1)~Every number, place name and event in the output can be traced to a ledger fact.
(P2)~The language model decides what to say and how to say it, but never supplies a value; values come from the ledger.
(P3)~Uncertainty estimated in Layers~1--2 reaches the text as calibrated hedges.
(P4)~When verification fails, a deterministic realizer takes over, so the SLM can add coverage but cannot remove faithfulness.

\begin{figure}[t]
  \centering
  \fbox{\parbox{0.95\linewidth}{\centering\vspace{2em}Architecture figure placeholder: IMU+GNSS $\rightarrow$ Layer 1 (orientation, ISO weighting, events) $\rightarrow$ Layer 2 (map matching, MI-PELT, HRNUs) $\rightarrow$ Layer 3 (ledger, KST capsule, plan--ground--realize, verifier, template fallback) $\rightarrow$ narrative / answers\vspace{2em}}}
  \caption{RoadScribe architecture.}
  \label{fig:arch}
\end{figure}

% ---------------------------------------------------------------------
\subsection{Layer 1: Kinetic Signal Processing}
\label{sec:layer1}

\paragraph{Notation.}
The IMU delivers specific force $\mathbf f^b[n]\in\R^3$ and angular rate $\boldsymbol\omega^b[n]\in\R^3$ in the phone body frame $b$ at rate $f_s$ ($T_s=1/f_s$).
GNSS delivers speed $v_k$ at times $t_k$ (1~Hz).
The vehicle frame $v$ follows ISO~8855~\cite{iso8855} ($x$ forward, $y$ left, $z$ up).
We estimate $\mathbf R_{vb}\in SO(3)$ and the gravity-free acceleration
\begin{equation}
  \mathbf a^v[n]=\mathbf R_{vb}\,\mathbf f^b[n]-g\,\mathbf e_z .
  \label{eq:av}
\end{equation}
Speed is interpolated to IMU time with a monotone cubic (PCHIP), giving $v(t)$ and $a^G_{\mathrm{lon}}(t)=\dot v(t)$.

\paragraph{Mount episodes.}
A phone in a holder is static relative to the vehicle except when handled.
We mark candidate handling where the gyroscope reports rotation that the vehicle does not perform, $\lVert\bar{\boldsymbol\omega}^b[n]\rVert-|\dot\psi^G(t_n)|>\Omega_h$, with $\dot\psi^G$ the GNSS course rate.
A candidate becomes an episode boundary when the quasi-static gravity directions on its two sides differ by more than $\theta_h$.
Orientation is estimated per episode $e$.

\paragraph{Gravity and tilt.}
Let $\bar{\mathbf f}^b$ and $\bar{\boldsymbol\omega}^b$ be low-passed at 0.5~Hz. The quasi-static set of episode $e$ is
\begin{equation}
\mathcal Q_e=\bigl\{n:\ \bigl|\lVert\bar{\mathbf f}^b[n]\rVert-g\bigr|<\epsilon_g,\ \lVert\bar{\boldsymbol\omega}^b[n]\rVert<\epsilon_\omega,\ |a^G_{\mathrm{lon}}(t_n)|<\epsilon_a,\ v(t_n)>v_{\min}\bigr\},
\end{equation}
and the up direction is $\mathbf u_e=\hat{\mathbf g}_e/\lVert\hat{\mathbf g}_e\rVert$ with $\hat{\mathbf g}_e$ the geometric median of $\{\bar{\mathbf f}^b[n]:n\in\mathcal Q_e\}$, which is robust to outliers and equivariant to rotation.
Gating matters: an ungated low-pass estimate absorbs sustained acceleration and tilts by $\arctan(a/g)$, i.e.\ $17^\circ$ at $3~\mathrm{m\,s^{-2}}$.
With $\mathbf k=\mathbf u_e\times\mathbf e_z/\lVert\mathbf u_e\times\mathbf e_z\rVert$ and $\theta=\operatorname{atan2}(\lVert\mathbf u_e\times\mathbf e_z\rVert,\mathbf u_e^\top\mathbf e_z)$, Rodrigues' formula gives the leveling rotation
\begin{equation}
\mathbf R_1=\mathbf I+\sin\theta\,[\mathbf k]_\times+(1-\cos\theta)\,[\mathbf k]_\times^2 .
\end{equation}
The route-average road grade biases this tilt (a 5\% grade is $2.9^\circ$).
The resulting error is at most 0.13\% on the vertical channel, but up to $0.5~\mathrm{m\,s^{-2}}$ on the longitudinal channel.
We therefore take longitudinal maneuver thresholds from GNSS rather than from the IMU.

\paragraph{Heading.}
After leveling, the horizontal axes differ from the vehicle axes by an unknown yaw $\psi$.
The leveled yaw rate $\omega^\ell_z=(\mathbf R_1\bar{\boldsymbol\omega}^b)_z$ does not depend on $\psi$.
This yields two GNSS-referenced targets at every fix: the longitudinal acceleration $a^G_{\mathrm{lon},k}$ and the centripetal acceleration $a^G_{\mathrm{lat},k}=v_k\,\omega^\ell_z(t_k)$.
With the leveled horizontal specific force $\mathbf h_k\in\R^2$ and $\mathbf r_k=(a^G_{\mathrm{lon},k},a^G_{\mathrm{lat},k})^\top$,
\begin{equation}
\hat\psi=\arg\min_{\psi}\sum_k\bigl\lVert\mathbf R_2(\psi)\mathbf h_k-\mathbf r_k\bigr\rVert^2
=\operatorname{atan2}\bigl(M_{21}-M_{12},\ M_{11}+M_{22}\bigr),\qquad \mathbf M=\sum_k\mathbf h_k\mathbf r_k^\top,
\label{eq:yaw}
\end{equation}
where $\mathbf R_2(\psi)=\bigl[\begin{smallmatrix}\cos\psi&\sin\psi\\-\sin\psi&\cos\psi\end{smallmatrix}\bigr]$.
The closed form follows from maximizing $\operatorname{tr}(\mathbf R_2\mathbf M)$.
Only fixes with $v_k>3~\mathrm{m\,s^{-1}}$ are used.
Finally $\mathbf R_{vb}=\mathbf R_z(\hat\psi)\mathbf R_1$, and a bootstrap over fixes provides the orientation confidence passed to Layer~3.

\paragraph{Speed decoupling.}
Vibration grows with speed.
For a road whose elevation PSD follows ISO~8608~\cite{iso8608}, $G_d(n)=G_d(n_0)(n/n_0)^{-w}$, the temporal input PSD at speed $v$ scales as $v^{w-1}f^{-w}$.
For a fixed vehicle response the RMS response therefore scales as $v^{(w-1)/2}$.
Rather than assume $w$, we estimate the elasticity with the window-level model $\log a_{w,i}=\alpha_{j(i)}+\gamma\log v_i+\varepsilon_i$, where $j$ indexes homogeneous segments, and report
\begin{equation}
\tilde a_w=a_w\,(v_{\mathrm{ref}}/\bar v)^{\hat\gamma},\qquad \bar v\ge v_{\min},\ v_{\mathrm{ref}}=10~\mathrm{m\,s^{-1}}.
\end{equation}
The raw $a_w$ describes how the ride felt and $\tilde a_w$ describes the road; narratives keep the two apart.

\paragraph{Frequency weighting.}
ISO~2631-1~\cite{iso2631} weights seated whole-body vibration with $W_k$ (vertical) and $W_d$ (horizontal).
Each is a product of band-limiting filters, an acceleration--velocity transition and, for $W_k$, an upward step ($p=j2\pi f$, $\omega_i=2\pi f_i$, $Q_1=1/\sqrt2$):
\begin{align}
H_h(p)&=\frac{p^2}{p^2+\omega_1p/Q_1+\omega_1^2}, &
H_l(p)&=\frac{\omega_2^2}{p^2+\omega_2p/Q_1+\omega_2^2}, \nonumber\\
H_t(p)&=\frac{(\omega_4^2/\omega_3)(p+\omega_3)}{p^2+\omega_4p/Q_4+\omega_4^2}, &
H_s(p)&=\frac{p^2+\omega_5p/Q_5+\omega_5^2}{p^2+\omega_6p/Q_6+\omega_6^2},
\end{align}
with $W=H_hH_lH_tH_s$ and the parameters of Table~\ref{tab:iso}.
Narratives are produced after the trip, so we apply $W$ in the frequency domain, $a_w=\mathcal F^{-1}\{\mathcal F\{a\}\,W(j2\pi f)\}$ with 10~s of zero padding.
This reproduces the analog magnitude and phase exactly below $f_s/2$.
For streaming use, a bilinear transform per second-order section, pre-warped at each section's corner frequency, is an alternative.
At $f_s=100$~Hz its magnitude error is at most 0.72~dB ($W_k$) and 1.22~dB ($W_d$) up to 20~Hz, but it reaches 7--8~dB at 40~Hz and underestimates impulsive peaks by 3.8\%.
At $f_s=400$~Hz the error stays within 0.23~dB up to 40~Hz.
The 100~Hz band limit $H_l$ cannot be realized for $f_s\le200$~Hz; omitting it changes $W$ by less than 0.11~dB below 40~Hz.
At $f_s=100$~Hz our analysis band ends near 40~Hz, below the 80~Hz upper limit of the standard.

\begin{table}[t]
\caption{ISO 2631-1 frequency-weighting parameters (Annex A). Frequencies in Hz.}
\label{tab:iso}
\begin{tabular}{lccccccccc}
\toprule
 & $f_1$ & $f_2$ & $f_3$ & $f_4$ & $Q_4$ & $f_5$ & $Q_5$ & $f_6$ & $Q_6$\\
\midrule
$W_k$ & 0.4 & 100 & 12.5 & 12.5 & 0.63 & 2.37 & 0.91 & 3.35 & 0.91\\
$W_d$ & 0.4 & 100 & 2.0 & 2.0 & 0.63 & -- & -- & -- & --\\
\bottomrule
\end{tabular}
\end{table}

\paragraph{Vibration measures.}
For each HRNU of duration $T$ and axis $j$ we compute
\begin{gather}
a_{w,j}=\Bigl(\tfrac1T\textstyle\int_0^Ta_{w,j}^2\,dt\Bigr)^{1/2},\qquad
a_v=\bigl(a_{w,x}^2+a_{w,y}^2+a_{w,z}^2\bigr)^{1/2},\\
\mathrm{MTVV}_j=\max_{t_0}\Bigl(\tfrac1\tau\textstyle\int_{t_0-\tau}^{t_0}a_{w,j}^2\,dt\Bigr)^{1/2}\ (\tau=1~\mathrm{s}),\qquad
\mathrm{VDV}_j=\Bigl(\textstyle\int_0^Ta_{w,j}^4\,dt\Bigr)^{1/4},\qquad
\mathrm{CF}_j=\frac{\max_t|a_{w,j}(t)|}{a_{w,j}}.
\end{gather}
Following the standard's evaluation rule, $a_w$ is always reported.
MTVV and VDV are added, and the unit is flagged as shock-dominated, when
\begin{equation}
S_j=\ind\bigl[\mathrm{CF}_j>9\ \lor\ \mathrm{MTVV}_j/a_{w,j}>1.5\ \lor\ \mathrm{VDV}_j/(a_{w,j}T^{1/4})>1.75\bigr].
\end{equation}
For a zero-mean stationary signal $\mathrm{VDV}/(a_wT^{1/4})=\kappa^{1/4}$, where $\kappa$ is the kurtosis.
The VDV criterion is therefore the condition $\kappa>9.38$; Gaussian vibration gives a ratio of 1.32.
Comfort descriptors follow the informative comfort scale of the standard (Annex~C).
Its bands overlap, so a single label uses geometric-mean cut-points $\boldsymbol\theta=(0.315,0.561,0.894,1.414,2.236)~\mathrm{m\,s^{-2}}$.
ISO~2631-1 defines comfort at the seat--occupant interface, whereas a phone in a holder measures the cabin structure.
We therefore call these quantities \emph{weighted cabin acceleration} and validate them against a phone fixed at the seat rail (Section~\ref{sec:data}).
Below about 1~Hz, $a_{w,x}$ and $a_{w,y}$ mostly reflect braking and cornering, and we report them as driving rather than road measures.

\paragraph{Modal band ratios.}
From the Welch PSD of $a^v_z$ (Hann, 5.12~s segments, 50\% overlap) we compute energies in a body band $[0.8,2.5]$~Hz, a wheel-hop band centred on the vehicle's wheel-hop peak $\hat f_{wh}\in[6,20]$~Hz, and a high band $[16,40]$~Hz.
From these we form $R_{wb}=\log(E_{\mathrm{wheel}}/E_{\mathrm{body}})$ and $R_{hb}=\log(E_{\mathrm{high}}/E_{\mathrm{body}})$.
Through $\lambda=v/f$ these bands correspond to wavelengths from about 0.25~m to 12.5~m at 10~m\,s$^{-1}$, i.e.\ unevenness and megatexture, not pavement macrotexture (0.5--50~mm~\cite{iso13473}).
They describe how the vehicle's modes are excited, and we use them within a vehicle and for segmentation.

\paragraph{Impulsive events and maneuver context.}
We band-pass the vertical channel, $\tilde z=\mathrm{BP}_{[4,25]\,\mathrm{Hz}}(a^v_z)$, and apply the Teager--Kaiser energy operator~\cite{kaiser1990tkeo},
\begin{equation}
\Psi[n]=\tilde z[n]^2-\tilde z[n-1]\,\tilde z[n+1],\qquad E[n]=\bigl(h_{50\,\mathrm{ms}}*\max(\Psi,0)\bigr)[n].
\end{equation}
For $\tilde z=A\cos(\Omega n+\phi)$, $\Psi=A^2\sin^2\Omega$, so the operator responds to both amplitude and frequency.
Candidates are maximal runs where $E$ exceeds $\theta[n]=\operatorname{med}_W(E)+\kappa\cdot1.4826\operatorname{MAD}_W(E)$ over a 10~s centred window.
The adaptive threshold follows the background, which is much higher on cobblestone than on asphalt.
Runs closer than 0.3~s are merged, and runs longer than $D_{\max}=0.5~\mathrm{s}+3~\mathrm{m}/\bar v$ are passed to segmentation as rough stretches.
The maneuver state $m_e$ is one of \textsc{brake}, \textsc{accel}, \textsc{turn}, \textsc{swerve}, \textsc{stop-go} or \textsc{cruise}.
Longitudinal states come from GNSS ($a^G_{\mathrm{lon}}$) and lateral ones from the yaw rate ($v\,\omega^v_z$); \textsc{stop-go} requires $v<2~\mathrm{m\,s^{-1}}$ and a longitudinal jerk $|j_x|>J_{sg}$.
With $\rho_e=E_{[8,25]}(e)/E_{[0.5,4]}(e)$ the band-energy ratio of $a^v_z$ in the event window, each candidate is labelled
\begin{equation}
\ell(e)=
\begin{cases}
\text{\textsc{artifact}} & m_e\in\{\text{\textsc{stop-go}},\text{\textsc{accel}}\}\ \land\ \max_{n\in e}|j_x[n]|>J_x\ \land\ \rho_e<\rho_0,\\
\text{\textsc{defect}} & \lnot\,\text{\textsc{artifact}}\ \land\ \rho_e\ge\rho_0\ \land\ d_e\le D_{\max},\\
\text{\textsc{ambiguous}} & \text{otherwise}.
\end{cases}
\label{eq:gate}
\end{equation}
Braking and turning never veto a defect.
In our data they do not increase vertical vibration (Section~\ref{sec:l1eval}), and drivers slow down \emph{before} speed bumps, so a braking veto would remove true events.
Acceleration is gated because it is the only maneuver that adds a dashboard-only high-band component.
A swerve without a vertical candidate within $\pm2$~s yields a separate low-confidence \emph{avoidance} event.
A logistic model trained on held-out routes maps the event features to $p_{\mathrm{defect}}$.

% ---------------------------------------------------------------------
\subsection{Layer 2: Spatial Grounding and Segmentation}
\label{sec:layer2}

\paragraph{Map matching.}
We match GNSS fixes to OpenStreetMap with a hidden Markov model~\cite{newson2009hmm}.
The states are candidate projections $x_{k,i}$ within 50~m of fix $z_k$, the emissions are $\mathcal N(\lVert z_k-x_{k,i}\rVert_{gc};0,\sigma_z^2)$, and the transitions are $\beta^{-1}\exp(-|\lVert z_k-z_{k+1}\rVert_{gc}-\lVert x_{k,i}-x_{k+1,j}\rVert_{\mathrm{route}}|/\beta)$, decoded with Viterbi.
Forward--backward marginals give a per-fix match confidence.
The matched arc length $s(t)$ replaces $\int v\,dt$.
Each matched way contributes its \texttt{surface}, \texttt{smoothness}, \texttt{highway} and \texttt{name} tags, and nodes such as \texttt{traffic\_calming} and level crossings become candidate anchors.

\paragraph{Map-informed change-point detection.}
We resample Layer-1 features on a 5~m along-track grid, $\mathbf y_j=[\log\tilde a_{w,z},\log a_{w,z},R_{wb},R_{hb},\text{event rate},\log a_{w,y}]_j$, robust-standardized per trip.
Speed is excluded so that boundaries do not follow speed changes.
Segmentation solves
\begin{equation}
\min_{m,\,\tau_{1:m}}\ \sum_{i=1}^{m+1}C\bigl(\mathbf y_{\tau_{i-1}+1:\tau_i}\bigr)+\sum_{i=1}^{m}\beta(\tau_i),\qquad
C(\mathbf y_{a:b})=\sum_{j=a}^{b}\lVert\mathbf y_j-\bar{\mathbf y}_{a:b}\rVert^2_{\mathbf D^{-1}},
\label{eq:pelt}
\end{equation}
with a penalty that falls near OpenStreetMap transitions $b\in\mathcal B$ (surface-tag changes, road-class changes, junctions):
\begin{equation}
\beta(\tau)=\beta_0\bigl(1-\lambda\,\pi(\tau)\bigr),\qquad
\pi(\tau)=\max_{b\in\mathcal B}w_b\exp\!\Bigl(-\frac{(s_\tau-s_b)^2}{2\sigma_b^2}\Bigr),\qquad \lambda\in[0,1).
\end{equation}
Setting $\lambda=0$ recovers PELT~\cite{killick2012pelt}.
Let $F(0)=0$, $G(t)=F(t)+\beta(t)\ind[t>0]$ and $F(s)=\min_{t\in\mathcal R_s}\{G(t)+C(\mathbf y_{t+1:s})\}$.

\begin{lemma}[Pruning under location-dependent penalties]
\label{lem:prune}
Suppose $C(\mathbf y_{t+1:s})+C(\mathbf y_{s+1:T})\le C(\mathbf y_{t+1:T})$ for all $t<s<T$.
If $G(t)+C(\mathbf y_{t+1:s})\ge G(s)$, then $t$ is not the optimal last change point for any $T>s$ and can be removed from $\mathcal R$.
\end{lemma}
\begin{proof}
For $T>s$, $G(t)+C(\mathbf y_{t+1:T})\ge G(t)+C(\mathbf y_{t+1:s})+C(\mathbf y_{s+1:T})\ge G(s)+C(\mathbf y_{s+1:T})$.
\end{proof}

The usual condition $F(t)+C(\mathbf y_{t+1:s})\ge F(s)$ coincides with Lemma~\ref{lem:prune} only for a constant penalty.
With a lower penalty at map transitions it discards valid candidates.
In 300 random series it returned a sub-optimal partition every time, whereas the rule of Lemma~\ref{lem:prune} matched exhaustive optimal partitioning in all 300.
We choose $\beta_0=c\,d\ln n$ and $(c,\lambda)$ by leave-one-route-out validation, and report the penalty path~\cite{haynes2017crops}.

\paragraph{Arbitration and HRNUs.}
A change point within 25~m of a relevant map transition is snapped to it: the map fixes the location and the signal decides existence.
A map transition without a change point is kept only if a two-sample Hotelling test on $\mathbf y$ across it rejects at $\alpha=0.01$; otherwise every way split in the map would start a new unit.
Units shorter than 50~m are merged into their closer neighbour.
Each boundary is typed as map+signal, signal-only or map-only.
Each unit carries a map--signal disagreement score $1-P(\text{surface}=\text{OSM tag}\mid\mathbf y)$ when the map has a surface tag.
The resulting HRNUs, with their measures, events, anchors and confidences, are the input to Layer~3.

% ---------------------------------------------------------------------
\subsection{Layer 3: Kinetic--Spatial Tokenization and Grounded Generation}
\label{sec:layer3}

\paragraph{Claim ledger.}
Layer~3 first converts HRNUs and events into a ledger $\mathcal L$ of typed facts with identifiers, values, tolerances and confidences.
The types are severity, numeric, event, anchor, comparison, missing and metadata.
A \textsc{missing} fact records that a value is unavailable and why, e.g.\ a 48~s GNSS gap.
The ledger is the reference for both the run-time verifier and the evaluation.

\paragraph{Kinetic--Spatial Tokenizer.}
Algorithm~\ref{alg:kst} compiles the ledger into a capsule that fits a token budget $B$.
Each unit has a salience $\sigma_m$ and three serializations (omit, brief, full) with token costs $c_{m,\ell}$ measured with the target tokenizer.
Selection solves a multiple-choice knapsack exactly by dynamic programming:
\begin{equation}
V_m(b)=\max_{\ell\in\Lambda_m}\bigl\{V_{m-1}(b-c_{m,\ell})+\sigma_m\eta_\ell\bigr\},\qquad V_0(b)=0,
\label{eq:knap}
\end{equation}
where $\Lambda_m$ excludes \textsc{omit} for units in the must-mention set $\mathcal M$ and $\eta_{\textsc{omit}}=0$.
The cost is $O(3MB)$ for $M$ units.

\begin{algorithm}[t]
\caption{Kinetic--Spatial Tokenizer (KST)}
\label{alg:kst}
\begin{algorithmic}[1]
\Require HRNUs $\mathcal U=(u_1,\dots,u_M)$ in route order; events $\mathcal E$; anchors $\mathcal A$; optional history $\mathcal H$; audience $\pi$; budget $B$; tokenizer $\mathrm{tok}$
\Ensure capsule $\kappa$ with $|\mathrm{tok}(\kappa)|\le B$; ledger $\mathcal L$
\State $\mathcal L\gets\textsc{MetaFacts}(\text{channels},\text{map coverage},\pi)$
\For{$m=1,\dots,M$}
  \State $k_m\gets\textsc{IsoClass}(a_{v,m})$;\quad $\tilde k_m\gets\textsc{IsoClass}(\tilde a_{w,m})$ \Comment{felt vs.\ road}
  \State $\mathcal E_m\gets\{e\in\mathcal E: s(e)\in[s^0_m,s^1_m)\}$
  \State $\alpha_m\gets$ nearest anchor in $\mathcal A$ within $r_{\max}$ of $u_m$, or $\varnothing$
  \State $\sigma_m\gets w_1k_m+w_2\sum_{e\in\mathcal E_m}p_e\,\mathrm{sev}(e)+w_3\,\Delta_{\mathcal H}(u_m)+w_4\,\delta^{\mathrm{map}}_m$
  \State $\mathcal L\gets\mathcal L\cup\textsc{Facts}(u_m,k_m,\tilde k_m,\mathcal E_m,\alpha_m)$ \Comment{including \textsc{missing} facts}
  \For{$\ell\in\{\textsc{omit},\textsc{brief},\textsc{full}\}$}
    \State $c_{m,\ell}\gets|\mathrm{tok}(\textsc{Ser}_\ell(u_m))|$
  \EndFor
\EndFor
\State $\mathcal M\gets\{m:\exists e\in\mathcal E_m,\ p_e\,\mathrm{sev}(e)\ge\eta^\star\}$ \Comment{must-mention units}
\State $B'\gets B-|\mathrm{tok}(\textsc{Header})|-c_{\mathrm{rest}}$
\State $\mathbf x^\star\gets\textsc{McKnapsack}(\{c_{m,\ell}\},\{\sigma_m\eta_\ell\},B',\mathcal M)$ \Comment{Eq.~\eqref{eq:knap}}
\If{$\mathbf x^\star$ is infeasible}
  \State merge adjacent units outside $\mathcal M$ and repeat; if still infeasible, increase $B$ and log it
\EndIf
\State $\kappa\gets\textsc{Header}\ \|\ \textsc{Ser}_{x^\star_1}(u_1)\ \|\cdots\|\ \textsc{Ser}_{x^\star_M}(u_M)\ \|\ \textsc{Rest}(\{u_m:x^\star_m=\textsc{omit}\})$
\State \Return $\kappa,\ \mathcal L$
\end{algorithmic}
\end{algorithm}

\paragraph{Plan--ground--realize.}
Generation runs in four steps.
\emph{Plan:} the SLM emits a JSON plan listing which units and facts to mention, their order, a rhetorical relation for each (elaboration, contrast, cause, comparison with history) and a hedge.
Decoding is constrained by a grammar generated per capsule from the ledger identifiers~\cite{ggerganov_llamacpp}, so the plan is always well-formed and cannot reference a unit or event that does not exist.
\emph{Ground:} deterministic code expands each plan item into a fact frame with values, names and units taken from $\mathcal L$.
\emph{Realize:} the SLM writes text from the fact frames, either with placeholders for numbers and names that are filled afterwards, or as free text.
\emph{Verify:} claims extracted from the text (Section~\ref{sec:metrics}) are checked against $\mathcal L$.
A failing sentence is regenerated at most twice and then replaced by the template realization of the same plan item.
The SLM therefore controls selection, ordering, aggregation and phrasing, but not values.

\paragraph{Hedging.}
The plan assigns hedges from fact confidence $p$: assertive for $p\ge0.9$, \emph{likely} for $0.7\le p<0.9$, \emph{possibly} for $0.4\le p<0.7$; facts with $p<0.4$ are mentioned only when asked about.
The verifier rejects a sentence whose hedge is stronger than its confidence band allows.

% =====================================================================
% Section 4 -- Evaluation
% Usage: % =====================================================================
% Section 4 -- Evaluation
% Usage: % =====================================================================
% Section 4 -- Evaluation
% Usage: \input{sec4_evaluation} after sec3_system (shares its macros).
% Bibliography keys used below (add to your .bib):
%   menegazzo2021pvs, min2023factscore, wiseman2017challenges, kent1964wep,
%   mielke2022linguistic, williams1949designs, judd2012stimuli, jian2000trust,
%   christensen_ordinal, wobbrock2011art, braun2006thematic, demontjoye2013unique,
%   thompson2026localising, vandeemter2005templates
% =====================================================================
\providecommand{\R}{\mathbb{R}}
\providecommand{\ind}{\mathbf{1}}
\providecommand{\TBD}[1][]{\textcolor{red}{[TBD#1]}}

\section{Evaluation}
\label{sec:eval}

We evaluate RoadScribe on five questions.
\begin{description}
\item[RQ1] Does Layer~1 recover vehicle-frame vibration and road events from a phone, and are its outputs grounded in vibration rather than in speed?
\item[RQ2] Does map-informed segmentation recover road-surface boundaries better than signal-only or map-only segmentation?
\item[RQ3] How faithful are the generated narratives to the ledger and to ground truth, and how do errors split between sensing and language?
\item[RQ4] Where does a grounded SLM add coverage over a template narrator, and at what fidelity cost?
\item[RQ5] How do people judge the narratives for accuracy, usefulness, trust and annoyance, and do they notice errors?
\end{description}

% ---------------------------------------------------------------------
\subsection{Data}
\label{sec:data}

\paragraph{The PVS dataset.}
The Passive Vehicular Sensors dataset~\cite{menegazzo2021pvs} contains nine trips that cover three routes (13.8, 11.6 and 10.7~km), each driven once by three vehicle--driver pairs.
In total there are 180.1~min and 108.6~km of driving.
IMU modules at the dashboard, above the suspension and below the suspension log at 100~Hz on each side of the vehicle, GNSS logs at 1~Hz, and an environment camera records each trip.
Labels mark the surface type (asphalt, cobblestone, dirt; 54 transitions) and speed bumps (60 traversals of 20 physical bumps).
We use the left-side logger: the dashboard unit serves as the phone proxy, and the suspension units serve as reference sensors that a phone does not have.
The left and right loggers are not synchronized.
Cross-correlating their yaw rates shows offsets of up to 0.54~s that drift within a trip, so we do not average the two sides.
PVS uses fixed modules, which exercise neither phone handling nor holder dynamics.

\paragraph{Phone corpus.}
To cover what PVS cannot, we recorded \TBD{} drives (\TBD{}~km, \TBD{} vehicles) with two phones.
One phone was in a dashboard holder and recorded video; the other was fixed to the seat rail as a floor reference.
Both logged at \TBD{}~Hz and were synchronized by tap events and yaw-rate cross-correlation (residual offset \TBD{}~ms).
The corpus includes \TBD{} repetitions of one route by one driver and \TBD{} drives with deliberate phone handling.
Two annotators marked speed breakers, potholes, rough patches, surface changes and rail crossings from video without seeing any narrative (Cohen's $\kappa=\TBD{}$ on \TBD{}\% double-coded).

\paragraph{Splits.}
All learned components, thresholds and prompts are fit with routes held out.
On PVS we use leave-one-route-out folds, and for vehicle generalization we test each (route, vehicle) trip using only the four trips that share neither its route nor its vehicle.
Evaluation questions, test personas and capsules from test drives never appear in prompts, exemplars or training data.

% ---------------------------------------------------------------------
\subsection{Narrators Compared}
\label{sec:conditions}

\begin{table}[t]
\caption{Narrators compared. All share Layers 1--2 and the ledger.}
\label{tab:conditions}
\begin{tabular}{lp{0.62\linewidth}}
\toprule
Code & Description \\
\midrule
T & Template NLG pipeline: rule-based content selection, aggregation, referring expressions, ISO-class lexical choice, history comparison, a fallback for every missing field (\TBD{} rules). \\
T+V & T with sampled lexical and ordering variation (controls for variety alone). \\
T$\rightarrow$P & T followed by SLM paraphrase, as in deployed driver-feedback systems~\cite{thompson2026localising}. \\
S-free & SLM given the capsule, free generation, no plan or verifier. \\
S-PGR & Plan--ground--realize with verification and template fallback, on device (ours). \\
C-PGR & S-PGR protocol with a cloud model (\TBD{}), as an upper reference. \\
\bottomrule
\end{tabular}
\end{table}

% ---------------------------------------------------------------------
\subsection{Metrics}
\label{sec:metrics}

\paragraph{Claims.}
A narrative $N$ is decomposed into atomic claims $\mathcal X(N)$ of type severity, numeric, event, anchor, comparison or hedge, in the spirit of~\cite{min2023factscore,wiseman2017challenges}.
The same LLM-based extractor is applied to every narrator.
On \TBD{} claims stratified by narrator it agreed with two annotators at $\kappa=\TBD{}$.
A claim is supported by reference $R$ when its value, unit, place and direction agree with $R$ within type-specific tolerances.
We use two references: the ledger $L$, which isolates the language layer, and ground truth $G$, which measures the system end to end:
\begin{equation}
P_R(N)=\frac{1}{|\mathcal X(N)|}\sum_{c\in\mathcal X(N)}\mathrm{supp}_R(c),\qquad R\in\{L,G\}.
\end{equation}
We report claims about facts marked \textsc{missing} separately; their rate should be zero.

\paragraph{Kinetic--semantic fidelity.}
Let $\mathcal S$ be the severity claims of a narrator pooled over the corpus.
Each claim $c\in\mathcal S$ pairs a measured intensity $x_c$ ($a_v$ for how the ride felt, $\tilde a_{w,z}$ for the road) with an expressed class $\hat k_c\in\{0,\dots,5\}$ from a fixed lexicon of ISO comfort descriptors.
With $R_c$ and $Q_c$ the mid-ranks of $x_c$ and $\hat k_c$,
\begin{equation}
\rho_{\mathrm{KSF}}=\frac{\sum_{c}(R_c-\bar R)(Q_c-\bar Q)}{\sqrt{\sum_{c}(R_c-\bar R)^2\sum_{c}(Q_c-\bar Q)^2}} .
\end{equation}
Because $\rho_{\mathrm{KSF}}$ ignores monotone distortion, we also report band consistency and an overstatement index,
\begin{equation}
\mathrm{BC}=\frac1{|\mathcal S|}\sum_{c\in\mathcal S}\ind\bigl[x_c\in\mathcal B(\hat k_c)\bigr],\qquad
\mathrm{OI}=\frac1{|\mathcal S|}\sum_{c\in\mathcal S}\bigl(\hat k_c-k^\star(x_c)\bigr),\qquad
k^\star(x)=\sum_{j=1}^5\ind[x\ge\theta_j],
\end{equation}
where $\mathcal B(k)$ is the (overlapping) ISO band of class $k$ and $\boldsymbol\theta$ are the cut-points of Section~\ref{sec:layer1}.

\paragraph{Anomaly grounding.}
Narrated event mentions $\hat{\mathcal E}$ are expanded by count; vague quantifiers expand to the lower bound of a fixed lexicon.
Each mention has a type and a resolved along-track interval $I(\hat e)$, and each reference event $e\in\mathcal E^R$ has a type and a position $s(e)$.
A mention and a reference event may match when their types agree and $\mathrm{dist}(s(e),I(\hat e))\le\delta_s$.
With $M^\star$ a maximum bipartite matching,
\begin{equation}
\mathrm{AGP}_R=\frac{|M^\star|}{|\hat{\mathcal E}|},\qquad
\mathrm{AGR}_R=\frac{|M^\star|}{|\mathcal E^R|},\qquad
\mathrm{wAGR}_R=\frac{\sum_{e\in M^\star}\mathrm{MTVV}_e}{\sum_{e\in\mathcal E^R}\mathrm{MTVV}_e}.
\end{equation}
For $R=L$, recall is computed over the must-mention events.
For $R=G$, the references are the labelled bumps (PVS) and the annotated events (phone corpus).
We use $\delta_s=25$~m, which covers label timing uncertainty and map-match error, and report 10 and 50~m in the supplement.
Crossing the two references attributes every error to perception (missed or false detection), selection (omitted) or verbalization (unsupported).

\paragraph{Calibration of hedges.}
For event claims with hedge $h\in\{\text{assertive},\text{likely},\text{possibly}\}$, let $P_G(h)$ be their precision against ground truth and $p_h$ the probability participants associate with the hedge, elicited in a short survey (cf.~\cite{kent1964wep,mielke2022linguistic}). Then
\begin{equation}
\mathrm{ECE}_{\mathrm{ling}}=\sum_h\frac{n_h}{n}\,\bigl|P_G(h)-p_h\bigr| .
\end{equation}

\paragraph{Novelty--change alignment.}
For consecutive drives $(d,d')$ of one route by one driver, let $\Delta_{\mathrm{kin}}$ be the mean absolute change of $\log a_v$ over 100~m bins plus the event-set Jaccard distance.
Let $\Delta_{\mathrm{claim}}$ be the Jaccard distance between canonicalized claim sets.
Then $\mathrm{NCA}=\mathrm{Spearman}(\Delta_{\mathrm{kin}},\Delta_{\mathrm{claim}})$ over pairs.
Lexical variation independent of content, $\mathrm{LV}=1-\mathrm{BLEU}$ between narratives with numbers and names masked, is reported alongside; variety that does not track change raises LV but not NCA.

\paragraph{Question answering and cost.}
For questions we report the correct-answer rate (judged against the ledger by raters blind to the narrator), the unsupported-claim rate, and abstention precision and recall on unanswerable questions.
System cost comprises time to first token, decode rate, end-to-end latency (p50/p95), peak memory, energy per narrative, plan validity, verifier rejection rate and template fallback rate.

% ---------------------------------------------------------------------
\subsection{Layer 1 and Layer 2}
\label{sec:l1eval}

\paragraph{Orientation.}
On PVS, the heading estimate of Eq.~\eqref{eq:yaw} recovered the known mount: $87.9^\circ$--$92.5^\circ$ for two vehicles (nominal $90^\circ$).
For the third vehicle it gave $103.1^\circ$--$105.5^\circ$ consistently across its three routes, indicating a mounting offset of about $14^\circ$.
Under 20 random three-dimensional rotations of the input per trip, the estimate rotated with the data to within $0.63^\circ$.
On the phone corpus, handling detection and re-estimation achieve \TBD{}.

\paragraph{Do maneuvers inflate vertical vibration?}
Table~\ref{tab:maneuver} compares two-second windows on asphalt (no speed bump within $\pm2$~s, 5--20~m\,s$^{-1}$) across maneuver states, controlling for speed and trip.
Braking lowered dashboard vertical vibration and turning had no detectable effect.
Acceleration added high-band energy at the dashboard that the below-suspension sensor did not see, consistent with powertrain vibration.
This is why the gate of Eq.~\eqref{eq:gate} targets acceleration and stop-and-go events and never vetoes on braking.
A wheelbase-lag cue (front then rear axle) did not separate road events: the lag between the two largest impacts did not scale with $1/v$ (Spearman $|\rho|\le0.12$, $n=60$), so it is not used.

\begin{table}[t]
\caption{Difference in log RMS vertical acceleration relative to cruising (PVS, asphalt, 2-s windows; OLS with log speed and trip effects, trip-clustered standard errors). Cruise $n=225$.}
\label{tab:maneuver}
\begin{tabular}{lcccc}
\toprule
Maneuver ($n$) & Dash 0.5--4 Hz & Dash 4--25 Hz & Wheel 4--25 Hz & Dash $-$ wheel, 4--25 Hz\\
\midrule
Brake (75) & $-0.18^{***}$ & $-0.13^{*}$ & $-0.16$ & $+0.03$\\
Accelerate (23) & $-0.18$ & $-0.06$ & $-0.37^{**}$ & $+0.31^{***}$\\
Turn (92) & $+0.07$ & $-0.04$ & $+0.06$ & $-0.10^{***}$\\
\bottomrule
\multicolumn{5}{l}{\footnotesize $^{*}p<.05$, $^{**}p<.01$, $^{***}p<.001$. Speed elasticity $\hat\gamma$: 0.32, 0.36, 0.62 (columns 1--3).}
\end{tabular}
\end{table}

\paragraph{Events and segmentation.}
Bump detection on held-out routes reaches recall \TBD{} at \TBD{} false alarms per km.
For comparison, a speed-only detector reaches \TBD{} and $\mathrm{RMS}(a_z)/\sqrt v$ reaches \TBD{}.
The gate changes false alarms by \TBD{} at fixed recall.
Against the 54 surface transitions, MI-PELT reaches boundary F1 \TBD{} at 25~m, compared with \TBD{} for constant-penalty PELT, \TBD{} for map-only boundaries and \TBD{} for PELT on speed alone.
OpenStreetMap carries a surface tag on \TBD{}\% of the matched length and a \texttt{traffic\_calming} node at \TBD{} of the 20 bumps.
Results on PVS cover three route clusters and are descriptive; inferential comparisons use the phone corpus.

% ---------------------------------------------------------------------
\subsection{Narrative Benchmark}
\label{sec:bench}

\paragraph{Tasks.}
Five task families probe where a template narrator suffices and where it does not.
\textbf{E1} answers questions from a sealed half of \TBD{} questions elicited from drivers before the system was built; the template baseline uses an intent classifier over intents designed on the other half.
\textbf{E2} writes for audiences that no template or exemplar targets (road engineer, motion-sensitive passenger, fleet manager); the template baseline uses the nearest designed audience.
\textbf{E3} narrates degraded or contradictory capsules: no gyroscope, a 60~s GNSS gap, failed map matching, map--signal disagreement, low-confidence events.
\textbf{E4} narrates repeated drives of one route.
\textbf{E5} grows the template narrator rule by rule on the development split and traces its coverage against the number of rules.
Thresholds that would count against our hypotheses were fixed before the sealed sets were used (\TBD[: registration statement]).

\paragraph{Results.}
\TBD[: Table with $P_L$, $P_G$, hallucinated-missing rate, $\rho_{\mathrm{KSF}}$, BC, OI, AGP/AGR ($L$ and $G$), $\mathrm{ECE}_{\mathrm{ling}}$, NCA, LV, QA correctness, fallback rate per narrator; frontier plot of $P_L$ against task coverage; error attribution table.]

% ---------------------------------------------------------------------
\subsection{On-Device Performance}
\label{sec:device}
\TBD[: two phones (one mid-range); model, quantization, time to first token, decode rate, end-to-end latency p50/p95, peak memory, energy per narrative, battery drain of 100--400~Hz logging, throttling over 10 consecutive generations. State which layers run on the phone.]

% ---------------------------------------------------------------------
\subsection{User Study}
\label{sec:study}

\paragraph{Design.}
Twenty-eight licensed drivers each judged narratives from four narrators (T, T+V, T$\rightarrow$P, S-PGR) in a within-subject design.
Narrator order followed a Williams design~\cite{williams1949designs} with sequences (T, T+V, S-PGR, T$\rightarrow$P), (T+V, T$\rightarrow$P, T, S-PGR), (T$\rightarrow$P, S-PGR, T+V, T) and (S-PGR, T, T$\rightarrow$P, T+V), each used by seven participants, which balances first-order carryover.
Every participant saw eight video excerpts of 2--3~min, two per narrator.
A second Latin square rotated excerpts across narrators, so each excerpt appeared with each narrator equally often and excerpts could be modelled as random effects~\cite{judd2012stimuli}.
In two of each participant's eight narratives one fact was altered (a count, a place or a severity), balanced across narrators, to measure whether errors were noticed.
Participants were debriefed about the altered narratives.

\paragraph{Participants and ethics.}
\TBD[: recruitment, inclusion criteria (licensed, $\ge$1 year of regular driving), demographics, compensation, ethics approval number]. Location traces are identifying~\cite{demontjoye2013unique}; they were stored encrypted and are released only truncated by 500~m at each end. No narrative was presented while a participant was driving.

\paragraph{Procedure.}
After consent and a practice trial, participants completed four blocks, one per narrator.
For each excerpt they watched the video, read the narrative and answered the items of Table~\ref{tab:items}.
They then answered whether anything seemed wrong, and if so, what.
Each block ended with the trust checklist of Jian et al.~\cite{jian2000trust}.
In a question-answering task, participants typed three questions about each of two further excerpts and rated the answers, which came from T with an intent classifier or from S-PGR, counterbalanced.
The session closed with a ranking of the narrators and a semi-structured interview, analysed thematically~\cite{braun2006thematic}.

\begin{table}[t]
\caption{Per-narrative items (7-point scale, 1 = strongly disagree, 7 = strongly agree).}
\label{tab:items}
\begin{tabular}{lp{0.78\linewidth}}
\toprule
Code & Item \\
\midrule
ACC1 & The summary accurately described the road conditions in the clip. \\
ACC2 & The places the summary mentioned matched where things happened in the clip. \\
SPEC & The summary was specific enough that I could find the places it mentioned again. \\
USE & This summary would help me decide how to drive this route next time. \\
CLAR & The summary was easy to understand on the first reading. \\
CERT & When the summary was unsure about something, it said so in a way that felt right. \\
NOVEL & The summary told me something about the drive that I had not noticed myself. \\
ANNOY & I would find getting a summary like this after every drive annoying. (reversed) \\
FIT & This summary was written for someone like me. \\
\bottomrule
\end{tabular}
\end{table}

\paragraph{Hypotheses.}
Four hypotheses were fixed before data collection.
H1: claim precision of S-PGR is non-inferior to T, $P_L(\text{S-PGR})-P_L(\text{T})>-0.02$.
H2: S-PGR answers the sealed questions correctly more often than T with an intent classifier.
H3: S-PGR is rated more useful (USE) than T+V.
H4: perceived accuracy, the mean of ACC1 and ACC2, of S-PGR is non-inferior to T within 0.5 scale points.
All other comparisons are exploratory.

\paragraph{Analysis.}
Ratings were analysed with cumulative link mixed models~\cite{christensen_ordinal},
\begin{equation}
\operatorname{logit}P(Y_{ij}\le k)=\theta_k-\bigl(\beta_{\mathrm{cond}(ij)}+\beta_{\mathrm{period}(ij)}+u_i+w_j\bigr),\qquad u_i\sim\mathcal N(0,\sigma_u^2),\ w_j\sim\mathcal N(0,\sigma_w^2),
\end{equation}
for participant $i$ and excerpt $j$, with Holm correction across H2 and H3 and a first-order carryover term as a check.
We also ran aligned-rank-transform ANOVA~\cite{wobbrock2011art} and Wilcoxon signed-rank tests on participant means as robustness checks.
H4 was tested on participant means of unaltered trials and holds if the lower bound of the 90\% confidence interval of the paired difference exceeds $-0.5$.
H1 used a cluster bootstrap over drives.
H2 and error detection used logistic mixed models with random intercepts for question or participant and for excerpt.

\paragraph{Power.}
With 28 participants, a paired comparison has 80\% power for standardized effects $d_z\ge0.55$ at $\alpha=.05$ ($d_z\ge0.65$ at the Holm-adjusted $\alpha=.05/3$).
The non-inferiority test of H4 has 82.5\% power when the margin equals half the standard deviation of paired differences.
The pilot ($n=\TBD{}$) estimated that standard deviation at \TBD{} scale points.

\paragraph{Field deployment.}
\TBD{} participants used RoadScribe on their own phones for 14 days.
Narrators alternated by day between T+V and S-PGR on a randomized balanced schedule.
Narratives were delivered only after the vehicle had been stationary for two minutes.
We logged opens, time to open, dwell time and questions asked, and collected three daily ratings (USE, ANNOY, NOVEL).
We model habituation as the narrator-by-day interaction in a mixed model with random slopes; with this sample size the analysis is descriptive.
 after sec3_system (shares its macros).
% Bibliography keys used below (add to your .bib):
%   menegazzo2021pvs, min2023factscore, wiseman2017challenges, kent1964wep,
%   mielke2022linguistic, williams1949designs, judd2012stimuli, jian2000trust,
%   christensen_ordinal, wobbrock2011art, braun2006thematic, demontjoye2013unique,
%   thompson2026localising, vandeemter2005templates
% =====================================================================
\providecommand{\R}{\mathbb{R}}
\providecommand{\ind}{\mathbf{1}}
\providecommand{\TBD}[1][]{\textcolor{red}{[TBD#1]}}

\section{Evaluation}
\label{sec:eval}

We evaluate RoadScribe on five questions.
\begin{description}
\item[RQ1] Does Layer~1 recover vehicle-frame vibration and road events from a phone, and are its outputs grounded in vibration rather than in speed?
\item[RQ2] Does map-informed segmentation recover road-surface boundaries better than signal-only or map-only segmentation?
\item[RQ3] How faithful are the generated narratives to the ledger and to ground truth, and how do errors split between sensing and language?
\item[RQ4] Where does a grounded SLM add coverage over a template narrator, and at what fidelity cost?
\item[RQ5] How do people judge the narratives for accuracy, usefulness, trust and annoyance, and do they notice errors?
\end{description}

% ---------------------------------------------------------------------
\subsection{Data}
\label{sec:data}

\paragraph{The PVS dataset.}
The Passive Vehicular Sensors dataset~\cite{menegazzo2021pvs} contains nine trips that cover three routes (13.8, 11.6 and 10.7~km), each driven once by three vehicle--driver pairs.
In total there are 180.1~min and 108.6~km of driving.
IMU modules at the dashboard, above the suspension and below the suspension log at 100~Hz on each side of the vehicle, GNSS logs at 1~Hz, and an environment camera records each trip.
Labels mark the surface type (asphalt, cobblestone, dirt; 54 transitions) and speed bumps (60 traversals of 20 physical bumps).
We use the left-side logger: the dashboard unit serves as the phone proxy, and the suspension units serve as reference sensors that a phone does not have.
The left and right loggers are not synchronized.
Cross-correlating their yaw rates shows offsets of up to 0.54~s that drift within a trip, so we do not average the two sides.
PVS uses fixed modules, which exercise neither phone handling nor holder dynamics.

\paragraph{Phone corpus.}
To cover what PVS cannot, we recorded \TBD{} drives (\TBD{}~km, \TBD{} vehicles) with two phones.
One phone was in a dashboard holder and recorded video; the other was fixed to the seat rail as a floor reference.
Both logged at \TBD{}~Hz and were synchronized by tap events and yaw-rate cross-correlation (residual offset \TBD{}~ms).
The corpus includes \TBD{} repetitions of one route by one driver and \TBD{} drives with deliberate phone handling.
Two annotators marked speed breakers, potholes, rough patches, surface changes and rail crossings from video without seeing any narrative (Cohen's $\kappa=\TBD{}$ on \TBD{}\% double-coded).

\paragraph{Splits.}
All learned components, thresholds and prompts are fit with routes held out.
On PVS we use leave-one-route-out folds, and for vehicle generalization we test each (route, vehicle) trip using only the four trips that share neither its route nor its vehicle.
Evaluation questions, test personas and capsules from test drives never appear in prompts, exemplars or training data.

% ---------------------------------------------------------------------
\subsection{Narrators Compared}
\label{sec:conditions}

\begin{table}[t]
\caption{Narrators compared. All share Layers 1--2 and the ledger.}
\label{tab:conditions}
\begin{tabular}{lp{0.62\linewidth}}
\toprule
Code & Description \\
\midrule
T & Template NLG pipeline: rule-based content selection, aggregation, referring expressions, ISO-class lexical choice, history comparison, a fallback for every missing field (\TBD{} rules). \\
T+V & T with sampled lexical and ordering variation (controls for variety alone). \\
T$\rightarrow$P & T followed by SLM paraphrase, as in deployed driver-feedback systems~\cite{thompson2026localising}. \\
S-free & SLM given the capsule, free generation, no plan or verifier. \\
S-PGR & Plan--ground--realize with verification and template fallback, on device (ours). \\
C-PGR & S-PGR protocol with a cloud model (\TBD{}), as an upper reference. \\
\bottomrule
\end{tabular}
\end{table}

% ---------------------------------------------------------------------
\subsection{Metrics}
\label{sec:metrics}

\paragraph{Claims.}
A narrative $N$ is decomposed into atomic claims $\mathcal X(N)$ of type severity, numeric, event, anchor, comparison or hedge, in the spirit of~\cite{min2023factscore,wiseman2017challenges}.
The same LLM-based extractor is applied to every narrator.
On \TBD{} claims stratified by narrator it agreed with two annotators at $\kappa=\TBD{}$.
A claim is supported by reference $R$ when its value, unit, place and direction agree with $R$ within type-specific tolerances.
We use two references: the ledger $L$, which isolates the language layer, and ground truth $G$, which measures the system end to end:
\begin{equation}
P_R(N)=\frac{1}{|\mathcal X(N)|}\sum_{c\in\mathcal X(N)}\mathrm{supp}_R(c),\qquad R\in\{L,G\}.
\end{equation}
We report claims about facts marked \textsc{missing} separately; their rate should be zero.

\paragraph{Kinetic--semantic fidelity.}
Let $\mathcal S$ be the severity claims of a narrator pooled over the corpus.
Each claim $c\in\mathcal S$ pairs a measured intensity $x_c$ ($a_v$ for how the ride felt, $\tilde a_{w,z}$ for the road) with an expressed class $\hat k_c\in\{0,\dots,5\}$ from a fixed lexicon of ISO comfort descriptors.
With $R_c$ and $Q_c$ the mid-ranks of $x_c$ and $\hat k_c$,
\begin{equation}
\rho_{\mathrm{KSF}}=\frac{\sum_{c}(R_c-\bar R)(Q_c-\bar Q)}{\sqrt{\sum_{c}(R_c-\bar R)^2\sum_{c}(Q_c-\bar Q)^2}} .
\end{equation}
Because $\rho_{\mathrm{KSF}}$ ignores monotone distortion, we also report band consistency and an overstatement index,
\begin{equation}
\mathrm{BC}=\frac1{|\mathcal S|}\sum_{c\in\mathcal S}\ind\bigl[x_c\in\mathcal B(\hat k_c)\bigr],\qquad
\mathrm{OI}=\frac1{|\mathcal S|}\sum_{c\in\mathcal S}\bigl(\hat k_c-k^\star(x_c)\bigr),\qquad
k^\star(x)=\sum_{j=1}^5\ind[x\ge\theta_j],
\end{equation}
where $\mathcal B(k)$ is the (overlapping) ISO band of class $k$ and $\boldsymbol\theta$ are the cut-points of Section~\ref{sec:layer1}.

\paragraph{Anomaly grounding.}
Narrated event mentions $\hat{\mathcal E}$ are expanded by count; vague quantifiers expand to the lower bound of a fixed lexicon.
Each mention has a type and a resolved along-track interval $I(\hat e)$, and each reference event $e\in\mathcal E^R$ has a type and a position $s(e)$.
A mention and a reference event may match when their types agree and $\mathrm{dist}(s(e),I(\hat e))\le\delta_s$.
With $M^\star$ a maximum bipartite matching,
\begin{equation}
\mathrm{AGP}_R=\frac{|M^\star|}{|\hat{\mathcal E}|},\qquad
\mathrm{AGR}_R=\frac{|M^\star|}{|\mathcal E^R|},\qquad
\mathrm{wAGR}_R=\frac{\sum_{e\in M^\star}\mathrm{MTVV}_e}{\sum_{e\in\mathcal E^R}\mathrm{MTVV}_e}.
\end{equation}
For $R=L$, recall is computed over the must-mention events.
For $R=G$, the references are the labelled bumps (PVS) and the annotated events (phone corpus).
We use $\delta_s=25$~m, which covers label timing uncertainty and map-match error, and report 10 and 50~m in the supplement.
Crossing the two references attributes every error to perception (missed or false detection), selection (omitted) or verbalization (unsupported).

\paragraph{Calibration of hedges.}
For event claims with hedge $h\in\{\text{assertive},\text{likely},\text{possibly}\}$, let $P_G(h)$ be their precision against ground truth and $p_h$ the probability participants associate with the hedge, elicited in a short survey (cf.~\cite{kent1964wep,mielke2022linguistic}). Then
\begin{equation}
\mathrm{ECE}_{\mathrm{ling}}=\sum_h\frac{n_h}{n}\,\bigl|P_G(h)-p_h\bigr| .
\end{equation}

\paragraph{Novelty--change alignment.}
For consecutive drives $(d,d')$ of one route by one driver, let $\Delta_{\mathrm{kin}}$ be the mean absolute change of $\log a_v$ over 100~m bins plus the event-set Jaccard distance.
Let $\Delta_{\mathrm{claim}}$ be the Jaccard distance between canonicalized claim sets.
Then $\mathrm{NCA}=\mathrm{Spearman}(\Delta_{\mathrm{kin}},\Delta_{\mathrm{claim}})$ over pairs.
Lexical variation independent of content, $\mathrm{LV}=1-\mathrm{BLEU}$ between narratives with numbers and names masked, is reported alongside; variety that does not track change raises LV but not NCA.

\paragraph{Question answering and cost.}
For questions we report the correct-answer rate (judged against the ledger by raters blind to the narrator), the unsupported-claim rate, and abstention precision and recall on unanswerable questions.
System cost comprises time to first token, decode rate, end-to-end latency (p50/p95), peak memory, energy per narrative, plan validity, verifier rejection rate and template fallback rate.

% ---------------------------------------------------------------------
\subsection{Layer 1 and Layer 2}
\label{sec:l1eval}

\paragraph{Orientation.}
On PVS, the heading estimate of Eq.~\eqref{eq:yaw} recovered the known mount: $87.9^\circ$--$92.5^\circ$ for two vehicles (nominal $90^\circ$).
For the third vehicle it gave $103.1^\circ$--$105.5^\circ$ consistently across its three routes, indicating a mounting offset of about $14^\circ$.
Under 20 random three-dimensional rotations of the input per trip, the estimate rotated with the data to within $0.63^\circ$.
On the phone corpus, handling detection and re-estimation achieve \TBD{}.

\paragraph{Do maneuvers inflate vertical vibration?}
Table~\ref{tab:maneuver} compares two-second windows on asphalt (no speed bump within $\pm2$~s, 5--20~m\,s$^{-1}$) across maneuver states, controlling for speed and trip.
Braking lowered dashboard vertical vibration and turning had no detectable effect.
Acceleration added high-band energy at the dashboard that the below-suspension sensor did not see, consistent with powertrain vibration.
This is why the gate of Eq.~\eqref{eq:gate} targets acceleration and stop-and-go events and never vetoes on braking.
A wheelbase-lag cue (front then rear axle) did not separate road events: the lag between the two largest impacts did not scale with $1/v$ (Spearman $|\rho|\le0.12$, $n=60$), so it is not used.

\begin{table}[t]
\caption{Difference in log RMS vertical acceleration relative to cruising (PVS, asphalt, 2-s windows; OLS with log speed and trip effects, trip-clustered standard errors). Cruise $n=225$.}
\label{tab:maneuver}
\begin{tabular}{lcccc}
\toprule
Maneuver ($n$) & Dash 0.5--4 Hz & Dash 4--25 Hz & Wheel 4--25 Hz & Dash $-$ wheel, 4--25 Hz\\
\midrule
Brake (75) & $-0.18^{***}$ & $-0.13^{*}$ & $-0.16$ & $+0.03$\\
Accelerate (23) & $-0.18$ & $-0.06$ & $-0.37^{**}$ & $+0.31^{***}$\\
Turn (92) & $+0.07$ & $-0.04$ & $+0.06$ & $-0.10^{***}$\\
\bottomrule
\multicolumn{5}{l}{\footnotesize $^{*}p<.05$, $^{**}p<.01$, $^{***}p<.001$. Speed elasticity $\hat\gamma$: 0.32, 0.36, 0.62 (columns 1--3).}
\end{tabular}
\end{table}

\paragraph{Events and segmentation.}
Bump detection on held-out routes reaches recall \TBD{} at \TBD{} false alarms per km.
For comparison, a speed-only detector reaches \TBD{} and $\mathrm{RMS}(a_z)/\sqrt v$ reaches \TBD{}.
The gate changes false alarms by \TBD{} at fixed recall.
Against the 54 surface transitions, MI-PELT reaches boundary F1 \TBD{} at 25~m, compared with \TBD{} for constant-penalty PELT, \TBD{} for map-only boundaries and \TBD{} for PELT on speed alone.
OpenStreetMap carries a surface tag on \TBD{}\% of the matched length and a \texttt{traffic\_calming} node at \TBD{} of the 20 bumps.
Results on PVS cover three route clusters and are descriptive; inferential comparisons use the phone corpus.

% ---------------------------------------------------------------------
\subsection{Narrative Benchmark}
\label{sec:bench}

\paragraph{Tasks.}
Five task families probe where a template narrator suffices and where it does not.
\textbf{E1} answers questions from a sealed half of \TBD{} questions elicited from drivers before the system was built; the template baseline uses an intent classifier over intents designed on the other half.
\textbf{E2} writes for audiences that no template or exemplar targets (road engineer, motion-sensitive passenger, fleet manager); the template baseline uses the nearest designed audience.
\textbf{E3} narrates degraded or contradictory capsules: no gyroscope, a 60~s GNSS gap, failed map matching, map--signal disagreement, low-confidence events.
\textbf{E4} narrates repeated drives of one route.
\textbf{E5} grows the template narrator rule by rule on the development split and traces its coverage against the number of rules.
Thresholds that would count against our hypotheses were fixed before the sealed sets were used (\TBD[: registration statement]).

\paragraph{Results.}
\TBD[: Table with $P_L$, $P_G$, hallucinated-missing rate, $\rho_{\mathrm{KSF}}$, BC, OI, AGP/AGR ($L$ and $G$), $\mathrm{ECE}_{\mathrm{ling}}$, NCA, LV, QA correctness, fallback rate per narrator; frontier plot of $P_L$ against task coverage; error attribution table.]

% ---------------------------------------------------------------------
\subsection{On-Device Performance}
\label{sec:device}
\TBD[: two phones (one mid-range); model, quantization, time to first token, decode rate, end-to-end latency p50/p95, peak memory, energy per narrative, battery drain of 100--400~Hz logging, throttling over 10 consecutive generations. State which layers run on the phone.]

% ---------------------------------------------------------------------
\subsection{User Study}
\label{sec:study}

\paragraph{Design.}
Twenty-eight licensed drivers each judged narratives from four narrators (T, T+V, T$\rightarrow$P, S-PGR) in a within-subject design.
Narrator order followed a Williams design~\cite{williams1949designs} with sequences (T, T+V, S-PGR, T$\rightarrow$P), (T+V, T$\rightarrow$P, T, S-PGR), (T$\rightarrow$P, S-PGR, T+V, T) and (S-PGR, T, T$\rightarrow$P, T+V), each used by seven participants, which balances first-order carryover.
Every participant saw eight video excerpts of 2--3~min, two per narrator.
A second Latin square rotated excerpts across narrators, so each excerpt appeared with each narrator equally often and excerpts could be modelled as random effects~\cite{judd2012stimuli}.
In two of each participant's eight narratives one fact was altered (a count, a place or a severity), balanced across narrators, to measure whether errors were noticed.
Participants were debriefed about the altered narratives.

\paragraph{Participants and ethics.}
\TBD[: recruitment, inclusion criteria (licensed, $\ge$1 year of regular driving), demographics, compensation, ethics approval number]. Location traces are identifying~\cite{demontjoye2013unique}; they were stored encrypted and are released only truncated by 500~m at each end. No narrative was presented while a participant was driving.

\paragraph{Procedure.}
After consent and a practice trial, participants completed four blocks, one per narrator.
For each excerpt they watched the video, read the narrative and answered the items of Table~\ref{tab:items}.
They then answered whether anything seemed wrong, and if so, what.
Each block ended with the trust checklist of Jian et al.~\cite{jian2000trust}.
In a question-answering task, participants typed three questions about each of two further excerpts and rated the answers, which came from T with an intent classifier or from S-PGR, counterbalanced.
The session closed with a ranking of the narrators and a semi-structured interview, analysed thematically~\cite{braun2006thematic}.

\begin{table}[t]
\caption{Per-narrative items (7-point scale, 1 = strongly disagree, 7 = strongly agree).}
\label{tab:items}
\begin{tabular}{lp{0.78\linewidth}}
\toprule
Code & Item \\
\midrule
ACC1 & The summary accurately described the road conditions in the clip. \\
ACC2 & The places the summary mentioned matched where things happened in the clip. \\
SPEC & The summary was specific enough that I could find the places it mentioned again. \\
USE & This summary would help me decide how to drive this route next time. \\
CLAR & The summary was easy to understand on the first reading. \\
CERT & When the summary was unsure about something, it said so in a way that felt right. \\
NOVEL & The summary told me something about the drive that I had not noticed myself. \\
ANNOY & I would find getting a summary like this after every drive annoying. (reversed) \\
FIT & This summary was written for someone like me. \\
\bottomrule
\end{tabular}
\end{table}

\paragraph{Hypotheses.}
Four hypotheses were fixed before data collection.
H1: claim precision of S-PGR is non-inferior to T, $P_L(\text{S-PGR})-P_L(\text{T})>-0.02$.
H2: S-PGR answers the sealed questions correctly more often than T with an intent classifier.
H3: S-PGR is rated more useful (USE) than T+V.
H4: perceived accuracy, the mean of ACC1 and ACC2, of S-PGR is non-inferior to T within 0.5 scale points.
All other comparisons are exploratory.

\paragraph{Analysis.}
Ratings were analysed with cumulative link mixed models~\cite{christensen_ordinal},
\begin{equation}
\operatorname{logit}P(Y_{ij}\le k)=\theta_k-\bigl(\beta_{\mathrm{cond}(ij)}+\beta_{\mathrm{period}(ij)}+u_i+w_j\bigr),\qquad u_i\sim\mathcal N(0,\sigma_u^2),\ w_j\sim\mathcal N(0,\sigma_w^2),
\end{equation}
for participant $i$ and excerpt $j$, with Holm correction across H2 and H3 and a first-order carryover term as a check.
We also ran aligned-rank-transform ANOVA~\cite{wobbrock2011art} and Wilcoxon signed-rank tests on participant means as robustness checks.
H4 was tested on participant means of unaltered trials and holds if the lower bound of the 90\% confidence interval of the paired difference exceeds $-0.5$.
H1 used a cluster bootstrap over drives.
H2 and error detection used logistic mixed models with random intercepts for question or participant and for excerpt.

\paragraph{Power.}
With 28 participants, a paired comparison has 80\% power for standardized effects $d_z\ge0.55$ at $\alpha=.05$ ($d_z\ge0.65$ at the Holm-adjusted $\alpha=.05/3$).
The non-inferiority test of H4 has 82.5\% power when the margin equals half the standard deviation of paired differences.
The pilot ($n=\TBD{}$) estimated that standard deviation at \TBD{} scale points.

\paragraph{Field deployment.}
\TBD{} participants used RoadScribe on their own phones for 14 days.
Narrators alternated by day between T+V and S-PGR on a randomized balanced schedule.
Narratives were delivered only after the vehicle had been stationary for two minutes.
We logged opens, time to open, dwell time and questions asked, and collected three daily ratings (USE, ANNOY, NOVEL).
We model habituation as the narrator-by-day interaction in a mixed model with random slopes; with this sample size the analysis is descriptive.
 after sec3_system (shares its macros).
% Bibliography keys used below (add to your .bib):
%   menegazzo2021pvs, min2023factscore, wiseman2017challenges, kent1964wep,
%   mielke2022linguistic, williams1949designs, judd2012stimuli, jian2000trust,
%   christensen_ordinal, wobbrock2011art, braun2006thematic, demontjoye2013unique,
%   thompson2026localising, vandeemter2005templates
% =====================================================================
\providecommand{\R}{\mathbb{R}}
\providecommand{\ind}{\mathbf{1}}
\providecommand{\TBD}[1][]{\textcolor{red}{[TBD#1]}}

\section{Evaluation}
\label{sec:eval}

We evaluate RoadScribe on five questions.
\begin{description}
\item[RQ1] Does Layer~1 recover vehicle-frame vibration and road events from a phone, and are its outputs grounded in vibration rather than in speed?
\item[RQ2] Does map-informed segmentation recover road-surface boundaries better than signal-only or map-only segmentation?
\item[RQ3] How faithful are the generated narratives to the ledger and to ground truth, and how do errors split between sensing and language?
\item[RQ4] Where does a grounded SLM add coverage over a template narrator, and at what fidelity cost?
\item[RQ5] How do people judge the narratives for accuracy, usefulness, trust and annoyance, and do they notice errors?
\end{description}

% ---------------------------------------------------------------------
\subsection{Data}
\label{sec:data}

\paragraph{The PVS dataset.}
The Passive Vehicular Sensors dataset~\cite{menegazzo2021pvs} contains nine trips that cover three routes (13.8, 11.6 and 10.7~km), each driven once by three vehicle--driver pairs.
In total there are 180.1~min and 108.6~km of driving.
IMU modules at the dashboard, above the suspension and below the suspension log at 100~Hz on each side of the vehicle, GNSS logs at 1~Hz, and an environment camera records each trip.
Labels mark the surface type (asphalt, cobblestone, dirt; 54 transitions) and speed bumps (60 traversals of 20 physical bumps).
We use the left-side logger: the dashboard unit serves as the phone proxy, and the suspension units serve as reference sensors that a phone does not have.
The left and right loggers are not synchronized.
Cross-correlating their yaw rates shows offsets of up to 0.54~s that drift within a trip, so we do not average the two sides.
PVS uses fixed modules, which exercise neither phone handling nor holder dynamics.

\paragraph{Phone corpus.}
To cover what PVS cannot, we recorded \TBD{} drives (\TBD{}~km, \TBD{} vehicles) with two phones.
One phone was in a dashboard holder and recorded video; the other was fixed to the seat rail as a floor reference.
Both logged at \TBD{}~Hz and were synchronized by tap events and yaw-rate cross-correlation (residual offset \TBD{}~ms).
The corpus includes \TBD{} repetitions of one route by one driver and \TBD{} drives with deliberate phone handling.
Two annotators marked speed breakers, potholes, rough patches, surface changes and rail crossings from video without seeing any narrative (Cohen's $\kappa=\TBD{}$ on \TBD{}\% double-coded).

\paragraph{Splits.}
All learned components, thresholds and prompts are fit with routes held out.
On PVS we use leave-one-route-out folds, and for vehicle generalization we test each (route, vehicle) trip using only the four trips that share neither its route nor its vehicle.
Evaluation questions, test personas and capsules from test drives never appear in prompts, exemplars or training data.

% ---------------------------------------------------------------------
\subsection{Narrators Compared}
\label{sec:conditions}

\begin{table}[t]
\caption{Narrators compared. All share Layers 1--2 and the ledger.}
\label{tab:conditions}
\begin{tabular}{lp{0.62\linewidth}}
\toprule
Code & Description \\
\midrule
T & Template NLG pipeline: rule-based content selection, aggregation, referring expressions, ISO-class lexical choice, history comparison, a fallback for every missing field (\TBD{} rules). \\
T+V & T with sampled lexical and ordering variation (controls for variety alone). \\
T$\rightarrow$P & T followed by SLM paraphrase, as in deployed driver-feedback systems~\cite{thompson2026localising}. \\
S-free & SLM given the capsule, free generation, no plan or verifier. \\
S-PGR & Plan--ground--realize with verification and template fallback, on device (ours). \\
C-PGR & S-PGR protocol with a cloud model (\TBD{}), as an upper reference. \\
\bottomrule
\end{tabular}
\end{table}

% ---------------------------------------------------------------------
\subsection{Metrics}
\label{sec:metrics}

\paragraph{Claims.}
A narrative $N$ is decomposed into atomic claims $\mathcal X(N)$ of type severity, numeric, event, anchor, comparison or hedge, in the spirit of~\cite{min2023factscore,wiseman2017challenges}.
The same LLM-based extractor is applied to every narrator.
On \TBD{} claims stratified by narrator it agreed with two annotators at $\kappa=\TBD{}$.
A claim is supported by reference $R$ when its value, unit, place and direction agree with $R$ within type-specific tolerances.
We use two references: the ledger $L$, which isolates the language layer, and ground truth $G$, which measures the system end to end:
\begin{equation}
P_R(N)=\frac{1}{|\mathcal X(N)|}\sum_{c\in\mathcal X(N)}\mathrm{supp}_R(c),\qquad R\in\{L,G\}.
\end{equation}
We report claims about facts marked \textsc{missing} separately; their rate should be zero.

\paragraph{Kinetic--semantic fidelity.}
Let $\mathcal S$ be the severity claims of a narrator pooled over the corpus.
Each claim $c\in\mathcal S$ pairs a measured intensity $x_c$ ($a_v$ for how the ride felt, $\tilde a_{w,z}$ for the road) with an expressed class $\hat k_c\in\{0,\dots,5\}$ from a fixed lexicon of ISO comfort descriptors.
With $R_c$ and $Q_c$ the mid-ranks of $x_c$ and $\hat k_c$,
\begin{equation}
\rho_{\mathrm{KSF}}=\frac{\sum_{c}(R_c-\bar R)(Q_c-\bar Q)}{\sqrt{\sum_{c}(R_c-\bar R)^2\sum_{c}(Q_c-\bar Q)^2}} .
\end{equation}
Because $\rho_{\mathrm{KSF}}$ ignores monotone distortion, we also report band consistency and an overstatement index,
\begin{equation}
\mathrm{BC}=\frac1{|\mathcal S|}\sum_{c\in\mathcal S}\ind\bigl[x_c\in\mathcal B(\hat k_c)\bigr],\qquad
\mathrm{OI}=\frac1{|\mathcal S|}\sum_{c\in\mathcal S}\bigl(\hat k_c-k^\star(x_c)\bigr),\qquad
k^\star(x)=\sum_{j=1}^5\ind[x\ge\theta_j],
\end{equation}
where $\mathcal B(k)$ is the (overlapping) ISO band of class $k$ and $\boldsymbol\theta$ are the cut-points of Section~\ref{sec:layer1}.

\paragraph{Anomaly grounding.}
Narrated event mentions $\hat{\mathcal E}$ are expanded by count; vague quantifiers expand to the lower bound of a fixed lexicon.
Each mention has a type and a resolved along-track interval $I(\hat e)$, and each reference event $e\in\mathcal E^R$ has a type and a position $s(e)$.
A mention and a reference event may match when their types agree and $\mathrm{dist}(s(e),I(\hat e))\le\delta_s$.
With $M^\star$ a maximum bipartite matching,
\begin{equation}
\mathrm{AGP}_R=\frac{|M^\star|}{|\hat{\mathcal E}|},\qquad
\mathrm{AGR}_R=\frac{|M^\star|}{|\mathcal E^R|},\qquad
\mathrm{wAGR}_R=\frac{\sum_{e\in M^\star}\mathrm{MTVV}_e}{\sum_{e\in\mathcal E^R}\mathrm{MTVV}_e}.
\end{equation}
For $R=L$, recall is computed over the must-mention events.
For $R=G$, the references are the labelled bumps (PVS) and the annotated events (phone corpus).
We use $\delta_s=25$~m, which covers label timing uncertainty and map-match error, and report 10 and 50~m in the supplement.
Crossing the two references attributes every error to perception (missed or false detection), selection (omitted) or verbalization (unsupported).

\paragraph{Calibration of hedges.}
For event claims with hedge $h\in\{\text{assertive},\text{likely},\text{possibly}\}$, let $P_G(h)$ be their precision against ground truth and $p_h$ the probability participants associate with the hedge, elicited in a short survey (cf.~\cite{kent1964wep,mielke2022linguistic}). Then
\begin{equation}
\mathrm{ECE}_{\mathrm{ling}}=\sum_h\frac{n_h}{n}\,\bigl|P_G(h)-p_h\bigr| .
\end{equation}

\paragraph{Novelty--change alignment.}
For consecutive drives $(d,d')$ of one route by one driver, let $\Delta_{\mathrm{kin}}$ be the mean absolute change of $\log a_v$ over 100~m bins plus the event-set Jaccard distance.
Let $\Delta_{\mathrm{claim}}$ be the Jaccard distance between canonicalized claim sets.
Then $\mathrm{NCA}=\mathrm{Spearman}(\Delta_{\mathrm{kin}},\Delta_{\mathrm{claim}})$ over pairs.
Lexical variation independent of content, $\mathrm{LV}=1-\mathrm{BLEU}$ between narratives with numbers and names masked, is reported alongside; variety that does not track change raises LV but not NCA.

\paragraph{Question answering and cost.}
For questions we report the correct-answer rate (judged against the ledger by raters blind to the narrator), the unsupported-claim rate, and abstention precision and recall on unanswerable questions.
System cost comprises time to first token, decode rate, end-to-end latency (p50/p95), peak memory, energy per narrative, plan validity, verifier rejection rate and template fallback rate.

% ---------------------------------------------------------------------
\subsection{Layer 1 and Layer 2}
\label{sec:l1eval}

\paragraph{Orientation.}
On PVS, the heading estimate of Eq.~\eqref{eq:yaw} recovered the known mount: $87.9^\circ$--$92.5^\circ$ for two vehicles (nominal $90^\circ$).
For the third vehicle it gave $103.1^\circ$--$105.5^\circ$ consistently across its three routes, indicating a mounting offset of about $14^\circ$.
Under 20 random three-dimensional rotations of the input per trip, the estimate rotated with the data to within $0.63^\circ$.
On the phone corpus, handling detection and re-estimation achieve \TBD{}.

\paragraph{Do maneuvers inflate vertical vibration?}
Table~\ref{tab:maneuver} compares two-second windows on asphalt (no speed bump within $\pm2$~s, 5--20~m\,s$^{-1}$) across maneuver states, controlling for speed and trip.
Braking lowered dashboard vertical vibration and turning had no detectable effect.
Acceleration added high-band energy at the dashboard that the below-suspension sensor did not see, consistent with powertrain vibration.
This is why the gate of Eq.~\eqref{eq:gate} targets acceleration and stop-and-go events and never vetoes on braking.
A wheelbase-lag cue (front then rear axle) did not separate road events: the lag between the two largest impacts did not scale with $1/v$ (Spearman $|\rho|\le0.12$, $n=60$), so it is not used.

\begin{table}[t]
\caption{Difference in log RMS vertical acceleration relative to cruising (PVS, asphalt, 2-s windows; OLS with log speed and trip effects, trip-clustered standard errors). Cruise $n=225$.}
\label{tab:maneuver}
\begin{tabular}{lcccc}
\toprule
Maneuver ($n$) & Dash 0.5--4 Hz & Dash 4--25 Hz & Wheel 4--25 Hz & Dash $-$ wheel, 4--25 Hz\\
\midrule
Brake (75) & $-0.18^{***}$ & $-0.13^{*}$ & $-0.16$ & $+0.03$\\
Accelerate (23) & $-0.18$ & $-0.06$ & $-0.37^{**}$ & $+0.31^{***}$\\
Turn (92) & $+0.07$ & $-0.04$ & $+0.06$ & $-0.10^{***}$\\
\bottomrule
\multicolumn{5}{l}{\footnotesize $^{*}p<.05$, $^{**}p<.01$, $^{***}p<.001$. Speed elasticity $\hat\gamma$: 0.32, 0.36, 0.62 (columns 1--3).}
\end{tabular}
\end{table}

\paragraph{Events and segmentation.}
Bump detection on held-out routes reaches recall \TBD{} at \TBD{} false alarms per km.
For comparison, a speed-only detector reaches \TBD{} and $\mathrm{RMS}(a_z)/\sqrt v$ reaches \TBD{}.
The gate changes false alarms by \TBD{} at fixed recall.
Against the 54 surface transitions, MI-PELT reaches boundary F1 \TBD{} at 25~m, compared with \TBD{} for constant-penalty PELT, \TBD{} for map-only boundaries and \TBD{} for PELT on speed alone.
OpenStreetMap carries a surface tag on \TBD{}\% of the matched length and a \texttt{traffic\_calming} node at \TBD{} of the 20 bumps.
Results on PVS cover three route clusters and are descriptive; inferential comparisons use the phone corpus.

% ---------------------------------------------------------------------
\subsection{Narrative Benchmark}
\label{sec:bench}

\paragraph{Tasks.}
Five task families probe where a template narrator suffices and where it does not.
\textbf{E1} answers questions from a sealed half of \TBD{} questions elicited from drivers before the system was built; the template baseline uses an intent classifier over intents designed on the other half.
\textbf{E2} writes for audiences that no template or exemplar targets (road engineer, motion-sensitive passenger, fleet manager); the template baseline uses the nearest designed audience.
\textbf{E3} narrates degraded or contradictory capsules: no gyroscope, a 60~s GNSS gap, failed map matching, map--signal disagreement, low-confidence events.
\textbf{E4} narrates repeated drives of one route.
\textbf{E5} grows the template narrator rule by rule on the development split and traces its coverage against the number of rules.
Thresholds that would count against our hypotheses were fixed before the sealed sets were used (\TBD[: registration statement]).

\paragraph{Results.}
\TBD[: Table with $P_L$, $P_G$, hallucinated-missing rate, $\rho_{\mathrm{KSF}}$, BC, OI, AGP/AGR ($L$ and $G$), $\mathrm{ECE}_{\mathrm{ling}}$, NCA, LV, QA correctness, fallback rate per narrator; frontier plot of $P_L$ against task coverage; error attribution table.]

% ---------------------------------------------------------------------
\subsection{On-Device Performance}
\label{sec:device}
\TBD[: two phones (one mid-range); model, quantization, time to first token, decode rate, end-to-end latency p50/p95, peak memory, energy per narrative, battery drain of 100--400~Hz logging, throttling over 10 consecutive generations. State which layers run on the phone.]

% ---------------------------------------------------------------------
\subsection{User Study}
\label{sec:study}

\paragraph{Design.}
Twenty-eight licensed drivers each judged narratives from four narrators (T, T+V, T$\rightarrow$P, S-PGR) in a within-subject design.
Narrator order followed a Williams design~\cite{williams1949designs} with sequences (T, T+V, S-PGR, T$\rightarrow$P), (T+V, T$\rightarrow$P, T, S-PGR), (T$\rightarrow$P, S-PGR, T+V, T) and (S-PGR, T, T$\rightarrow$P, T+V), each used by seven participants, which balances first-order carryover.
Every participant saw eight video excerpts of 2--3~min, two per narrator.
A second Latin square rotated excerpts across narrators, so each excerpt appeared with each narrator equally often and excerpts could be modelled as random effects~\cite{judd2012stimuli}.
In two of each participant's eight narratives one fact was altered (a count, a place or a severity), balanced across narrators, to measure whether errors were noticed.
Participants were debriefed about the altered narratives.

\paragraph{Participants and ethics.}
\TBD[: recruitment, inclusion criteria (licensed, $\ge$1 year of regular driving), demographics, compensation, ethics approval number]. Location traces are identifying~\cite{demontjoye2013unique}; they were stored encrypted and are released only truncated by 500~m at each end. No narrative was presented while a participant was driving.

\paragraph{Procedure.}
After consent and a practice trial, participants completed four blocks, one per narrator.
For each excerpt they watched the video, read the narrative and answered the items of Table~\ref{tab:items}.
They then answered whether anything seemed wrong, and if so, what.
Each block ended with the trust checklist of Jian et al.~\cite{jian2000trust}.
In a question-answering task, participants typed three questions about each of two further excerpts and rated the answers, which came from T with an intent classifier or from S-PGR, counterbalanced.
The session closed with a ranking of the narrators and a semi-structured interview, analysed thematically~\cite{braun2006thematic}.

\begin{table}[t]
\caption{Per-narrative items (7-point scale, 1 = strongly disagree, 7 = strongly agree).}
\label{tab:items}
\begin{tabular}{lp{0.78\linewidth}}
\toprule
Code & Item \\
\midrule
ACC1 & The summary accurately described the road conditions in the clip. \\
ACC2 & The places the summary mentioned matched where things happened in the clip. \\
SPEC & The summary was specific enough that I could find the places it mentioned again. \\
USE & This summary would help me decide how to drive this route next time. \\
CLAR & The summary was easy to understand on the first reading. \\
CERT & When the summary was unsure about something, it said so in a way that felt right. \\
NOVEL & The summary told me something about the drive that I had not noticed myself. \\
ANNOY & I would find getting a summary like this after every drive annoying. (reversed) \\
FIT & This summary was written for someone like me. \\
\bottomrule
\end{tabular}
\end{table}

\paragraph{Hypotheses.}
Four hypotheses were fixed before data collection.
H1: claim precision of S-PGR is non-inferior to T, $P_L(\text{S-PGR})-P_L(\text{T})>-0.02$.
H2: S-PGR answers the sealed questions correctly more often than T with an intent classifier.
H3: S-PGR is rated more useful (USE) than T+V.
H4: perceived accuracy, the mean of ACC1 and ACC2, of S-PGR is non-inferior to T within 0.5 scale points.
All other comparisons are exploratory.

\paragraph{Analysis.}
Ratings were analysed with cumulative link mixed models~\cite{christensen_ordinal},
\begin{equation}
\operatorname{logit}P(Y_{ij}\le k)=\theta_k-\bigl(\beta_{\mathrm{cond}(ij)}+\beta_{\mathrm{period}(ij)}+u_i+w_j\bigr),\qquad u_i\sim\mathcal N(0,\sigma_u^2),\ w_j\sim\mathcal N(0,\sigma_w^2),
\end{equation}
for participant $i$ and excerpt $j$, with Holm correction across H2 and H3 and a first-order carryover term as a check.
We also ran aligned-rank-transform ANOVA~\cite{wobbrock2011art} and Wilcoxon signed-rank tests on participant means as robustness checks.
H4 was tested on participant means of unaltered trials and holds if the lower bound of the 90\% confidence interval of the paired difference exceeds $-0.5$.
H1 used a cluster bootstrap over drives.
H2 and error detection used logistic mixed models with random intercepts for question or participant and for excerpt.

\paragraph{Power.}
With 28 participants, a paired comparison has 80\% power for standardized effects $d_z\ge0.55$ at $\alpha=.05$ ($d_z\ge0.65$ at the Holm-adjusted $\alpha=.05/3$).
The non-inferiority test of H4 has 82.5\% power when the margin equals half the standard deviation of paired differences.
The pilot ($n=\TBD{}$) estimated that standard deviation at \TBD{} scale points.

\paragraph{Field deployment.}
\TBD{} participants used RoadScribe on their own phones for 14 days.
Narrators alternated by day between T+V and S-PGR on a randomized balanced schedule.
Narratives were delivered only after the vehicle had been stationary for two minutes.
We logged opens, time to open, dwell time and questions asked, and collected three daily ratings (USE, ANNOY, NOVEL).
We model habituation as the narrator-by-day interaction in a mixed model with random slopes; with this sample size the analysis is descriptive.

% =====================================================================
% Section 5 -- Discussion, limitations, ethics, conclusion
% Usage: % =====================================================================
% Section 5 -- Discussion, limitations, ethics, conclusion
% Usage: % =====================================================================
% Section 5 -- Discussion, limitations, ethics, conclusion
% Usage: \input{sec5_discussion} inside an acmart document (see main.tex).
% =====================================================================
\providecommand{\TBD}[1][]{\textcolor{red}{[TBD#1]}}

\section{Discussion}
\label{sec:discussion}

\paragraph{When is a template enough?}
\TBD[: interpret E1--E5 and H1--H4. State plainly where T matched S-PGR (expected: fixed-schema summaries) and where it did not, with the fidelity cost of each gain. If the fallback rate exceeds 30\% on the sealed sets, report it as the main limitation.]

\paragraph{What the guardrails do and do not guarantee.}
Plan--ground--realize removes two failure classes by construction: the plan cannot reference a unit or fact that does not exist, and in the placeholder variant the model cannot state a number.
It does not make text correct by construction.
The verifier catches what its extractor can parse; on the template lexicon it rejected 98.9\% of controlled errors (Section~\ref{sec:verifier}), and its misses (small numeric drifts within tolerance, mentions that can be matched to a different event) show where free text can slip through.
On free text, the guarantee is only as good as the claim extractor, which is why we validate it against human annotation and report hallucinated-missing claims separately.

\paragraph{Sensing errors are not language errors.}
A narrative can be faithful to the ledger and still wrong about the road.
Scoring every claim against both the ledger and ground truth, and attributing each ground-truth event to perception, selection or verbalization, keeps the two apart.
Without that split, an improvement in the language layer could hide a regression in sensing, or the reverse.

\subsection{Limitations}
\label{sec:limits}

\begin{itemize}
\item \emph{PVS is not phone data.} Its fixed modules validate the signal processing and the segmentation, but not phone handling, holder resonance or seat-versus-holder transfer; these rest on our own two-phone corpus, which is smaller and from one region (\TBD{}).
\item \emph{Cabin, not seat.} ISO~2631-1 comfort is defined at the seat--occupant interface. Our measures describe the cabin structure; the seat-rail reference phone measures the seat's floor input, not the seat surface.
\item \emph{Bandwidth.} At 100~Hz the weighted band ends near 40~Hz, below the standard's 80~Hz.
\item \emph{Speed confound.} Surface type and speed are confounded in real driving. Speed decoupling with an estimated elasticity mitigates this, and every Layer~1 classifier is compared with speed-only baselines, but a residual confound cannot be excluded on three PVS routes.
\item \emph{Map coverage.} The map prior is only as good as OpenStreetMap's surface and traffic-calming tags on the driven roads; where coverage is low, the learned $\lambda$ approaches zero and MI-PELT reduces to PELT.
\item \emph{Synthetic data.} The synthetic corpus tests Layer~3 under controlled conditions; its narratives come from the template narrator and say nothing about what an SLM adds. Any SLM training on it is limited to the plan format, grounding and hedging, and excludes the E2 test audiences.
\item \emph{Sample size.} Twenty-eight participants support the four confirmatory hypotheses at the stated power, not fine-grained subgroup analyses; the field phase is descriptive.
\end{itemize}

\subsection{Ethics and Safety}
\label{sec:ethics}
Trip recordings reveal where people live and work~\cite{demontjoye2013unique}.
RoadScribe processes Layers~1 and~3 on the phone; traces used in research are stored encrypted and released only truncated or aggregated.
Narratives are delivered only after a trip has ended and the vehicle has been stationary, never while driving.
A system that narrates road conditions can be repurposed to monitor drivers, for example by insurers or employers; our narrators describe the road and the ride, attribute jolts to driving only to explain them, and are not designed or evaluated for scoring drivers.
The study \TBD[: ethics approval, consent, data retention].

\section{Conclusion}
\label{sec:conclusion}

RoadScribe turns phone vibration and map topology into grounded trip narratives.
Its sensing layer separates how a ride felt from how the road is, under the evaluation rules of ISO~2631-1.
Its segmentation uses the map without trusting it blindly, with a pruning rule that is exact for map-dependent penalties.
Its generation layer lets a small language model choose what to say while every value comes from a verified ledger.
\TBD[: one sentence with the headline benchmark and study results.]
We release the code, the tests that reproduce every verified number in this paper, and the synthetic benchmark data.
 inside an acmart document (see main.tex).
% =====================================================================
\providecommand{\TBD}[1][]{\textcolor{red}{[TBD#1]}}

\section{Discussion}
\label{sec:discussion}

\paragraph{When is a template enough?}
\TBD[: interpret E1--E5 and H1--H4. State plainly where T matched S-PGR (expected: fixed-schema summaries) and where it did not, with the fidelity cost of each gain. If the fallback rate exceeds 30\% on the sealed sets, report it as the main limitation.]

\paragraph{What the guardrails do and do not guarantee.}
Plan--ground--realize removes two failure classes by construction: the plan cannot reference a unit or fact that does not exist, and in the placeholder variant the model cannot state a number.
It does not make text correct by construction.
The verifier catches what its extractor can parse; on the template lexicon it rejected 98.9\% of controlled errors (Section~\ref{sec:verifier}), and its misses (small numeric drifts within tolerance, mentions that can be matched to a different event) show where free text can slip through.
On free text, the guarantee is only as good as the claim extractor, which is why we validate it against human annotation and report hallucinated-missing claims separately.

\paragraph{Sensing errors are not language errors.}
A narrative can be faithful to the ledger and still wrong about the road.
Scoring every claim against both the ledger and ground truth, and attributing each ground-truth event to perception, selection or verbalization, keeps the two apart.
Without that split, an improvement in the language layer could hide a regression in sensing, or the reverse.

\subsection{Limitations}
\label{sec:limits}

\begin{itemize}
\item \emph{PVS is not phone data.} Its fixed modules validate the signal processing and the segmentation, but not phone handling, holder resonance or seat-versus-holder transfer; these rest on our own two-phone corpus, which is smaller and from one region (\TBD{}).
\item \emph{Cabin, not seat.} ISO~2631-1 comfort is defined at the seat--occupant interface. Our measures describe the cabin structure; the seat-rail reference phone measures the seat's floor input, not the seat surface.
\item \emph{Bandwidth.} At 100~Hz the weighted band ends near 40~Hz, below the standard's 80~Hz.
\item \emph{Speed confound.} Surface type and speed are confounded in real driving. Speed decoupling with an estimated elasticity mitigates this, and every Layer~1 classifier is compared with speed-only baselines, but a residual confound cannot be excluded on three PVS routes.
\item \emph{Map coverage.} The map prior is only as good as OpenStreetMap's surface and traffic-calming tags on the driven roads; where coverage is low, the learned $\lambda$ approaches zero and MI-PELT reduces to PELT.
\item \emph{Synthetic data.} The synthetic corpus tests Layer~3 under controlled conditions; its narratives come from the template narrator and say nothing about what an SLM adds. Any SLM training on it is limited to the plan format, grounding and hedging, and excludes the E2 test audiences.
\item \emph{Sample size.} Twenty-eight participants support the four confirmatory hypotheses at the stated power, not fine-grained subgroup analyses; the field phase is descriptive.
\end{itemize}

\subsection{Ethics and Safety}
\label{sec:ethics}
Trip recordings reveal where people live and work~\cite{demontjoye2013unique}.
RoadScribe processes Layers~1 and~3 on the phone; traces used in research are stored encrypted and released only truncated or aggregated.
Narratives are delivered only after a trip has ended and the vehicle has been stationary, never while driving.
A system that narrates road conditions can be repurposed to monitor drivers, for example by insurers or employers; our narrators describe the road and the ride, attribute jolts to driving only to explain them, and are not designed or evaluated for scoring drivers.
The study \TBD[: ethics approval, consent, data retention].

\section{Conclusion}
\label{sec:conclusion}

RoadScribe turns phone vibration and map topology into grounded trip narratives.
Its sensing layer separates how a ride felt from how the road is, under the evaluation rules of ISO~2631-1.
Its segmentation uses the map without trusting it blindly, with a pruning rule that is exact for map-dependent penalties.
Its generation layer lets a small language model choose what to say while every value comes from a verified ledger.
\TBD[: one sentence with the headline benchmark and study results.]
We release the code, the tests that reproduce every verified number in this paper, and the synthetic benchmark data.
 inside an acmart document (see main.tex).
% =====================================================================
\providecommand{\TBD}[1][]{\textcolor{red}{[TBD#1]}}

\section{Discussion}
\label{sec:discussion}

\paragraph{When is a template enough?}
\TBD[: interpret E1--E5 and H1--H4. State plainly where T matched S-PGR (expected: fixed-schema summaries) and where it did not, with the fidelity cost of each gain. If the fallback rate exceeds 30\% on the sealed sets, report it as the main limitation.]

\paragraph{What the guardrails do and do not guarantee.}
Plan--ground--realize removes two failure classes by construction: the plan cannot reference a unit or fact that does not exist, and in the placeholder variant the model cannot state a number.
It does not make text correct by construction.
The verifier catches what its extractor can parse; on the template lexicon it rejected 98.9\% of controlled errors (Section~\ref{sec:verifier}), and its misses (small numeric drifts within tolerance, mentions that can be matched to a different event) show where free text can slip through.
On free text, the guarantee is only as good as the claim extractor, which is why we validate it against human annotation and report hallucinated-missing claims separately.

\paragraph{Sensing errors are not language errors.}
A narrative can be faithful to the ledger and still wrong about the road.
Scoring every claim against both the ledger and ground truth, and attributing each ground-truth event to perception, selection or verbalization, keeps the two apart.
Without that split, an improvement in the language layer could hide a regression in sensing, or the reverse.

\subsection{Limitations}
\label{sec:limits}

\begin{itemize}
\item \emph{PVS is not phone data.} Its fixed modules validate the signal processing and the segmentation, but not phone handling, holder resonance or seat-versus-holder transfer; these rest on our own two-phone corpus, which is smaller and from one region (\TBD{}).
\item \emph{Cabin, not seat.} ISO~2631-1 comfort is defined at the seat--occupant interface. Our measures describe the cabin structure; the seat-rail reference phone measures the seat's floor input, not the seat surface.
\item \emph{Bandwidth.} At 100~Hz the weighted band ends near 40~Hz, below the standard's 80~Hz.
\item \emph{Speed confound.} Surface type and speed are confounded in real driving. Speed decoupling with an estimated elasticity mitigates this, and every Layer~1 classifier is compared with speed-only baselines, but a residual confound cannot be excluded on three PVS routes.
\item \emph{Map coverage.} The map prior is only as good as OpenStreetMap's surface and traffic-calming tags on the driven roads; where coverage is low, the learned $\lambda$ approaches zero and MI-PELT reduces to PELT.
\item \emph{Synthetic data.} The synthetic corpus tests Layer~3 under controlled conditions; its narratives come from the template narrator and say nothing about what an SLM adds. Any SLM training on it is limited to the plan format, grounding and hedging, and excludes the E2 test audiences.
\item \emph{Sample size.} Twenty-eight participants support the four confirmatory hypotheses at the stated power, not fine-grained subgroup analyses; the field phase is descriptive.
\end{itemize}

\subsection{Ethics and Safety}
\label{sec:ethics}
Trip recordings reveal where people live and work~\cite{demontjoye2013unique}.
RoadScribe processes Layers~1 and~3 on the phone; traces used in research are stored encrypted and released only truncated or aggregated.
Narratives are delivered only after a trip has ended and the vehicle has been stationary, never while driving.
A system that narrates road conditions can be repurposed to monitor drivers, for example by insurers or employers; our narrators describe the road and the ride, attribute jolts to driving only to explain them, and are not designed or evaluated for scoring drivers.
The study \TBD[: ethics approval, consent, data retention].

\section{Conclusion}
\label{sec:conclusion}

RoadScribe turns phone vibration and map topology into grounded trip narratives.
Its sensing layer separates how a ride felt from how the road is, under the evaluation rules of ISO~2631-1.
Its segmentation uses the map without trusting it blindly, with a pruning rule that is exact for map-dependent penalties.
Its generation layer lets a small language model choose what to say while every value comes from a verified ledger.
\TBD[: one sentence with the headline benchmark and study results.]
We release the code, the tests that reproduce every verified number in this paper, and the synthetic benchmark data.


\begin{acks}
\TBD[: funding and acknowledgements; omitted for review.]
\end{acks}

\bibliographystyle{ACM-Reference-Format}
\bibliography{references}

\end{document}
